% !TEX root = ../../main.tex
\setchapterimage{figure/chapter6封面图.jpg}
\setchapterintro{}
\chapter{欧氏空间}
\section{内积及其性质}
\begin{definition}
	设 $V$ 是一个非零实线性空间. 若映射
	\[
	f:V\times V\longrightarrow \mathbb{R}
	\]
	满足下面三条性质，则称 $f$ 为 $V$ 上的一个\textbf{内积}，记作
	\[
	f(\alpha,\beta)=(\alpha,\beta)\quad \text{或}\quad \langle \alpha,\beta\rangle .
	\]
	\begin{enumerate}
		\item 双线性：对任意 $\alpha,\alpha_1,\alpha_2,\beta,\beta_1,\beta_2\in V$ 与 $k_1,k_2,l_1,l_2\in\mathbb{R}$，有
		\[
		\begin{aligned}
			(k_1\alpha_1+k_2\alpha_2,\beta)&=k_1(\alpha_1,\beta)+k_2(\alpha_2,\beta),\\
			(\alpha,l_1\beta_1+l_2\beta_2)&=l_1(\alpha,\beta_1)+l_2(\alpha,\beta_2).
		\end{aligned}
		\]
		\item 对称性：对任意 $\alpha,\beta\in V$，有
		\[
		(\alpha,\beta)=(\beta,\alpha).
		\]
		\item 正定性：对任意 $\alpha\in V$，有 $(\alpha,\alpha)\geqslant 0$，且
		\[
		(\alpha,\alpha)=0\Longleftrightarrow \alpha=0.
		\]
	\end{enumerate}
	配备内积的实线性空间 $(V,(\cdot,\cdot))$ 称为\textbf{实内积空间}. 若 $\dim V<\infty$，则称其为\textbf{欧氏空间}.
\end{definition}

\begin{example}
	在 $\mathbb{R}^n$ 中，取
	\[
	\alpha=\begin{pmatrix}x_1\\ \vdots\\ x_n\end{pmatrix},
	\qquad
	\beta=\begin{pmatrix}y_1\\ \vdots\\ y_n\end{pmatrix},
	\]
	定义
	\[
	(\alpha,\beta)=\alpha'\beta=\sum_{i=1}^{n}x_i y_i.
	\]
	这称为 $\mathbb{R}^n$ 上的标准内积.
\end{example}

\begin{proposition}
	每一个有限维实线性空间都可以赋予一个内积.
\end{proposition}
\begin{proof}
	设 $\dim V=n$，取 $V$ 的一组基 $\alpha_1,\cdots,\alpha_n$. 对任意 $\alpha,\beta\in V$，存在唯一的坐标向量 $X,Y\in\mathbb{R}^n$ 使得
	\[
	\alpha=(\alpha_1,\cdots,\alpha_n)X,
	\qquad
	\beta=(\alpha_1,\cdots,\alpha_n)Y.
	\]
	定义
	\[
	(\alpha,\beta)=X'Y.
	\]
	由 $\mathbb{R}^n$ 标准内积的双线性、对称性和正定性可知，上式给出 $V$ 上的一个内积.
\end{proof}

\begin{example}
	\begin{enumerate}
		\item 设 $w\in C[a,b]$ 且 $w(x)>0$，对任意 $f,g\in C[a,b]$，定义
		\[
		(f,g)=\int_a^b f(x)g(x)w(x)\,dx.
		\]
		则 $(\cdot,\cdot)$ 是 $C[a,b]$ 上的一个内积.
		\item 对任意 $f,g\in\mathbb{R}[x]$，定义
		\[
		(f,g)=\int_0^1 f(x)g(x)\,dx.
		\]
		则 $(\cdot,\cdot)$ 是 $\mathbb{R}[x]$ 上的一个内积.
	\end{enumerate}
	
\end{example}


\begin{lemma}
	设 $(V,(\cdot,\cdot))$ 是非零实内积空间，则有：
	\begin{enumerate}
		\item 对任意 $\alpha\in V$，有 $(\alpha,0)=0$.
		\item 若 $\beta\in V$ 且对任意 $\alpha\in V$ 都有 $(\alpha,\beta)=0$，则 $\beta=0$.
		\item 若 $W$ 是 $V$ 的子空间，则 $(\cdot,\cdot)$ 在 $W$ 上的限制仍然是 $W$ 上的内积.
	\end{enumerate}
\end{lemma}
\begin{proof}
	\begin{enumerate}
		\item 由双线性，$(\alpha,0)=(\alpha,0+0)=(\alpha,0)+(\alpha,0)$，故 $(\alpha,0)=0$.
		\item 令 $\alpha=\beta$，由假设得 $(\beta,\beta)=0$，再由正定性得 $\beta=0$.
		\item 双线性与对称性在子空间上显然保持；若 $\alpha\in W$ 且 $(\alpha,\alpha)=0$，则在 $V$ 中已有 $\alpha=0$. 因此限制仍是内积.
	\end{enumerate}
\end{proof}



\begin{definition}
	设 $(V,(\cdot,\cdot))$ 是实内积空间.
	\begin{enumerate}
		\item 若 $(\alpha,\beta)=0$，则称 $\alpha$ 与 $\beta$ \textbf{正交}，记作 $\alpha\perp\beta$.
		\item 向量 $\alpha$ 的\textbf{长度}或\textbf{范数}定义为
		\[
		|\alpha|=\sqrt{(\alpha,\alpha)}.
		\]
		也常记作 $\|\alpha\|$.
		\item 向量 $\alpha$ 与 $\beta$ 的\textbf{距离}定义为
		\[
		d(\alpha,\beta)=|\alpha-\beta|.
		\]
		\item 若 $\gamma\in V$ 且 $|\gamma|=1$，则称 $\gamma$ 为\textbf{单位向量}. 对 $\alpha\ne0$，向量
		\[
		\frac{1}{|\alpha|}\alpha
		\]
		是与 $\alpha$ 同方向的单位向量，这个过程称为\textbf{单位化}.
	\end{enumerate}
\end{definition}

\begin{theorem}(Cauchy--Bunyakovsky--Schwarz 不等式,简记为CBS不等式)
	设 $(V,(\cdot,\cdot))$ 是实内积空间. 对任意 $\alpha,\beta\in V$，有
	\[
	|(\alpha,\beta)|\leqslant |\alpha|\,|\beta|.
	\]
	等号成立当且仅当 $\alpha,\beta$ 线性相关.
\end{theorem}
\begin{proof}
	若 $\alpha=0$，结论显然成立. 下设 $\alpha\ne0$. 对任意 $t\in\mathbb{R}$，令
	\[
	\varphi(t)=(\beta+t\alpha,\beta+t\alpha).
	\]
	由正定性，$\varphi(t)\geqslant0$. 展开得
	\[
	\varphi(t)=(\alpha,\alpha)t^2+2(\alpha,\beta)t+(\beta,\beta).
	\]
	该二次多项式恒非负，故判别式不大于 $0$，即
	\[
	4(\alpha,\beta)^2-4(\alpha,\alpha)(\beta,\beta)\leqslant0.
	\]
	于是
	\[
	|(\alpha,\beta)|\leqslant |\alpha|\,|\beta|.
	\]
	若等号成立，则判别式为 $0$，从而存在 $t_0\in\mathbb{R}$ 使得
	\[
	0=\varphi(t_0)=(\beta+t_0\alpha,\beta+t_0\alpha),
	\]
	故 $\beta+t_0\alpha=0$，即 $\alpha,\beta$ 线性相关. 反之，若 $\alpha,\beta$ 线性相关，则可直接验证等号成立.
\end{proof}

由 Cauchy--Bunyakovsky--Schwarz 不等式可知，当 $\alpha,\beta\ne0$ 时，数
\[
\frac{(\alpha,\beta)}{|\alpha|\,|\beta|}
\]
属于区间 $[-1,1]$. 因此可以定义 $\alpha$ 与 $\beta$ 的夹角 $\theta$ 为满足
\[
\cos\theta=\frac{(\alpha,\beta)}{|\alpha|\,|\beta|}
\]
的角.

\begin{corollary}[三角不等式]
	对任意 $\alpha,\beta\in V$，有
	\[
	|\alpha+\beta|\leqslant |\alpha|+|\beta|.
	\]
\end{corollary}
\begin{proof}
	由 Cauchy--Bunyakovsky--Schwarz 不等式，
	\[
	\begin{aligned}
		|\alpha+\beta|^2
		&=(\alpha+\beta,\alpha+\beta)\\
		&=|\alpha|^2+2(\alpha,\beta)+|\beta|^2\\
		&\leqslant |\alpha|^2+2|\alpha|\,|\beta|+|\beta|^2\\
		&=(|\alpha|+|\beta|)^2.
	\end{aligned}
	\]
	两边开平方即得结论.
\end{proof}

\begin{example}
	设 $(V,(\cdot,\cdot))$ 是实内积空间，$\alpha_1,\cdots,\alpha_s\in V$ 两两正交，即
	\[
	(\alpha_i,\alpha_j)=0,
	\qquad 1\leqslant i\ne j\leqslant s.
	\]
	\begin{enumerate}
		\item 若 $\alpha_i\ne0$，$i=1,\cdots,s$，则 $\alpha_1,\cdots,\alpha_s$ 线性无关.
		\item 若 $\beta\in \Span\{\alpha_1,\cdots,\alpha_s\}$，则
		\[
		\beta=\frac{(\beta,\alpha_1)}{(\alpha_1,\alpha_1)}\alpha_1+\frac{(\beta,\alpha_2)}{(\alpha_2,\alpha_2)}\alpha_2+\cdots+
		\frac{(\beta,\alpha_s)}{(\alpha_s,\alpha_s)}\alpha_s.
		\]
	\end{enumerate}
\end{example}
\begin{proof}
	\begin{enumerate}
		\item 设 $t_1\alpha_1+\cdots+t_s\alpha_s=0$. 两边与 $\alpha_i$ 作内积，得
		\[
		t_i(\alpha_i,\alpha_i)=0.
		\]
		因 $\alpha_i\ne0$，故 $(\alpha_i,\alpha_i)>0$，于是 $t_i=0$. 对每个 $i$ 都成立，所以 $\alpha_1,\cdots,\alpha_s$ 线性无关.
		\item 设 $\beta=x_1\alpha_1+\cdots+x_s\alpha_s$. 两边与 $\alpha_i$ 作内积，得
		\[
		(\beta,\alpha_i)=x_i(\alpha_i,\alpha_i),
		\]
		故 $x_i=\dfrac{(\beta,\alpha_i)}{(\alpha_i,\alpha_i)}$.
	\end{enumerate}
\end{proof}
\section{正交性}
\begin{example}
	设 $(V,(\cdot,\cdot))$ 是实内积空间，固定 $0\ne\alpha\in V$. 定义
	\[
	t_\alpha:V\longrightarrow\mathbb{R},
	\qquad
	t_\alpha(\beta)=(\alpha,\beta).
	\]
	则 $t_\alpha\in\Hom(V,\mathbb{R})$，并且
	\[
	V=\ker t_\alpha\oplus \langle \alpha\rangle.
	\]
\end{example}
\begin{proof}
	由内积对第二个变量的线性性，$t_\alpha$ 是线性映射. 任取 $\beta\in V$，有
	\[
	\beta=\left(\beta-\frac{(\alpha,\beta)}{(\alpha,\alpha)}\alpha\right)
	+
	\frac{(\alpha,\beta)}{(\alpha,\alpha)}\alpha.
	\]
	并且
	\[
	t_\alpha\left(\beta-\frac{(\alpha,\beta)}{(\alpha,\alpha)}\alpha\right)
	=(\alpha,\beta)-\frac{(\alpha,\beta)}{(\alpha,\alpha)}(\alpha,\alpha)=0.
	\]
	故第一项属于 $\ker t_\alpha$，第二项属于 $\langle \alpha\rangle$. 若 $\gamma\in\ker t_\alpha\cap\langle \alpha\rangle$，设 $\gamma=k\alpha$，则
	\[
	0=t_\alpha(\gamma)=k(\alpha,\alpha),
	\]
	由于 $(\alpha,\alpha)>0$，得 $k=0$，故交为 $\{0\}$. 于是是直和.
\end{proof}


\begin{definition}
	设 $(V,(\cdot,\cdot))$ 是实内积空间，$0\ne\alpha\in V$，$\beta\in V$. 向量
	\[
	\operatorname{proj}_{\alpha}\beta=\frac{(\alpha,\beta)}{(\alpha,\alpha)}\alpha
	\]
	称为 $\beta$ 沿 $\alpha$ 的\textbf{正交投影}. 事实上，
	\[
	\left(\alpha,\beta-\frac{(\alpha,\beta)}{(\alpha,\alpha)}\alpha\right)=0,
	\]
	所以 $\beta-\operatorname{proj}_{\alpha}\beta$ 与 $\alpha$ 正交.
\end{definition}

\begin{definition}
	设 $(V,(\cdot,\cdot))$ 是实内积空间，$W$ 是 $V$ 的子空间. 定义
	\[
	W^{\perp}=\{\alpha\in V:(\alpha,\beta)=0,\ \forall\beta\in W\}.
	\]
	称 $W^{\perp}$ 为 $W$ 在 $V$ 中的\textbf{正交补}.
\end{definition}

\begin{proposition}
	设 $(V,(\cdot,\cdot))$ 是实内积空间，$W$ 是 $V$ 的子空间，则 $W^{\perp}$ 是 $V$ 的子空间，且
	\[
	W\cap W^{\perp}=\{0\}.
	\]
\end{proposition}
\begin{proof}
	若 $\alpha_1,\alpha_2\in W^{\perp}$，$k_1,k_2\in\mathbb{R}$，则对任意 $\beta\in W$，有
	\[
	(k_1\alpha_1+k_2\alpha_2,\beta)=k_1(\alpha_1,\beta)+k_2(\alpha_2,\beta)=0.
	\]
	故 $k_1\alpha_1+k_2\alpha_2\in W^{\perp}$，于是 $W^{\perp}$ 是子空间.
	若 $\alpha\in W\cap W^{\perp}$，则 $\alpha\in W$ 且 $(\alpha,\beta)=0$ 对任意 $\beta\in W$ 成立. 取 $\beta=\alpha$，得 $(\alpha,\alpha)=0$，故 $\alpha=0$. 因此 $W\cap W^{\perp}=\{0\}$.
\end{proof}


	一般的内积空间中，并不一定对所有子空间 $W$ 都有
	\[
	V=W\oplus W^{\perp}.
	\]
	在有限维情形中该结论总成立；在 Hilbert 空间中，通常需要 $W$ 为闭子空间.




\begin{theorem}(Schmidt 正交化)
	设 $(V,(\cdot,\cdot))$ 是实内积空间，$\alpha_1,\cdots,\alpha_m\in V$ 线性无关. 则存在 $\beta_1,\cdots,\beta_m\in V$，使得
	\begin{enumerate}
		\item $(\beta_i,\beta_j)=0$，$i\ne j$；
		\item 对任意 $1\leqslant k\leqslant m$，有$\alpha_1,\cdots,\alpha_k $与,$\beta_1,\cdots,\beta_k $等价
	\end{enumerate}
\end{theorem}
\begin{proof}
	构造如下：
	\[
	\begin{aligned}
		\beta_1&=\alpha_1,\\
		\beta_2&=\alpha_2-\frac{(\alpha_2,\beta_1)}{(\beta_1,\beta_1)}\beta_1,\\
		\beta_3&=\alpha_3-\frac{(\alpha_3,\beta_1)}{(\beta_1,\beta_1)}\beta_1
		-
		\frac{(\alpha_3,\beta_2)}{(\beta_2,\beta_2)}\beta_2,\\
		&\ \vdots\\
		\beta_k&=\alpha_k-\frac{(\alpha_k,\beta_1)}{(\beta_1,\beta_1)}\beta_1
		-\frac{(\alpha_k,\beta_2)}{(\beta_2,\beta_2)}\beta_2-\cdots-
		\frac{(\alpha_k,\beta_{k-1})}{(\beta_{k-1},\beta_{k-1})}\beta_{k-1}.
	\end{aligned}
	\]
	由于 $\alpha_1,\cdots,\alpha_m$ 线性无关，按归纳可知每个 $\beta_k\ne0$. 由构造式，对 $j<k$ 有
	\[
	(\beta_k,\beta_j)
	=(\alpha_k,\beta_j)-\frac{(\alpha_k,\beta_j)}{(\beta_j,\beta_j)}(\beta_j,\beta_j)=0.
	\]
	因此 $\beta_1,\cdots,\beta_m$ 两两正交. 同时，由构造式可知，每个 $\beta_k$ 属于 $\langle\alpha_1,\cdots,\alpha_k\rangle$，且
	\[
	\alpha_k\in\langle\beta_1,\cdots,\beta_k\rangle.
	\]
	故对每个 $k$ 都有
	\[
	\langle\alpha_1,\cdots,\alpha_k\rangle=\langle\beta_1,\cdots,\beta_k\rangle.
	\]


	也就是说，$\alpha_1,\cdots,\alpha_m$ 经过 Schmidt 正交化后得到 $\beta_1,\cdots,\beta_m$.
\end{proof}


	若把向量按列排列，则 Schmidt 正交化可以写为
		\[\left(\beta_1 \cdots \beta_m\right)=\left(\alpha_1 \cdots \alpha_m\right)\left(\begin{array}{ccc}
		1 & \cdots & * \\
		& \ddots & \vdots \\
		& & 1
	\end{array}\right) \]
	
	若继续令
	\[
	e_i=\frac{1}{|\beta_i|}\beta_i,
	\qquad i=1,\cdots,m,
	\]
	则 $e_1,\cdots,e_m$ 是一组标准正交向量.


\begin{theorem}
	设 $(V,(\cdot,\cdot))$ 是 $n$ 维实内积空间，$W$ 是 $V$ 的子空间，则
	\[
	V=W\oplus W^{\perp}.
	\]
\end{theorem}
\begin{proof}
	若 $W=\{0\}$，结论显然成立. 下设 $\dim W=r$，$0<r<n$. 取 $W$ 的一组基 $\alpha_1,\cdots,\alpha_r$，并扩充为 $V$ 的一组基
	\[
	\alpha_1,\cdots,\alpha_r,\alpha_{r+1},\cdots,\alpha_n.
	\]
	对这组基作 Schmidt 正交化，得到两两正交的非零向量
	\[
	\beta_1,\cdots,\beta_r,\beta_{r+1},\cdots,\beta_n.
	\]
	由 Schmidt 正交化的性质，
	\[
	W=\langle\beta_1,\cdots,\beta_r\rangle.
	\]
	令
	\[
	W_1=\langle\beta_{r+1},\cdots,\beta_n\rangle.
	\]
	显然 $V=W\oplus W_1$. 又因为 $\beta_{r+1},\cdots,\beta_n$ 与 $\beta_1,\cdots,\beta_r$ 正交，所以 $W_1\subset W^{\perp}$.
	反过来，任取 $\gamma\in W^{\perp}$. 由 $V=W\oplus W_1$，可唯一写成
	\[
	\gamma=w+w_1,
	\qquad w\in W,
	\quad w_1\in W_1.
	\]
	由于 $W_1\subset W^{\perp}$，有 $w_1\in W^{\perp}$. 因而 $w=\gamma-w_1\in W^{\perp}$. 但 $w\in W$，所以
	\[
	w\in W\cap W^{\perp}=\{0\}.
	\]
	故 $w=0$，从而 $\gamma=w_1\in W_1$. 于是 $W^{\perp}\subset W_1$，故 $W^{\perp}=W_1$. 因此
	\[
	V=W\oplus W^{\perp}.
	\]
\end{proof}
$ W^{\perp}$称为$W$的正交补空间
\begin{definition}
	在上一定理的条件下，对任意 $\alpha\in V$，存在唯一的分解
	\[
	\alpha=\alpha_1+\alpha_2,
	\qquad
	\alpha_1\in W,
	\quad
	\alpha_2\in W^{\perp}.
	\]
	称 $\alpha_1$ 为 $\alpha$ 在 $W$ 上的\textbf{正交投影}，称 $\alpha_2$ 为 $\alpha$ 到 $W$ 的\textbf{正交分量}.
\end{definition}


\begin{definition}
	设 $(V,(\cdot,\cdot))$ 是 $n$ 维欧氏空间. 若 $V$ 的一组基
	\[
	\varepsilon_1,\cdots,\varepsilon_n
	\]
	满足
	\[
	(\varepsilon_i,\varepsilon_j)=\delta_{ij},
	\qquad 1\leqslant i,j\leqslant n,
	\]
	其中
	\[
	\delta_{ij}=\begin{cases}
		1, & i=j,\\
		0, & i\ne j,
	\end{cases}
	\]
	则称 $\varepsilon_1,\cdots,\varepsilon_n$ 是 $V$ 的一组\textbf{标准正交基}，简称\textbf{标准基}.
\end{definition}

\begin{corollary}
	任何欧氏空间都有标准正交基.
\end{corollary}
\begin{proof}
	设 $\dim V=n$，取 $V$ 的一组基 $\alpha_1,\cdots,\alpha_n$. 对这组基作 Schmidt 正交化，得到两两正交的非零向量 $\beta_1,\cdots,\beta_n$. 再令
	\[
	e_i=\frac{1}{|\beta_i|}\beta_i,
	\qquad i=1,\cdots,n.
	\]
	则 $e_1,\cdots,e_n$ 是 $V$ 的一组标准正交基.
\end{proof}

\begin{corollary}
	设 $W$ 是 $V$ 的一个子空间. 对 $W$ 的一组标准正交基，可以延拓为 $V$ 的一组标准正交基.
\end{corollary}
\begin{proof}
	设 $\varepsilon_1,\cdots,\varepsilon_r$ 是 $W$ 的一组标准正交基. 若 $r=\dim V$，则结论显然成立. 若 $r<\dim V$，将其扩充为 $V$ 的一组基
	\[
	\varepsilon_1,\cdots,\varepsilon_r,\alpha_{r+1},\cdots,\alpha_n.
	\]
	对这组基作 Schmidt 正交化. 由于 $\varepsilon_1,\cdots,\varepsilon_r$ 已经是标准正交的，正交化并单位化后前 $r$ 个向量仍为 $\varepsilon_1,\cdots,\varepsilon_r$. 因此得到
	\[
	\varepsilon_1,\cdots,\varepsilon_r,e_{r+1},\cdots,e_n,
	\]
	这是 $V$ 的一组标准正交基.
\end{proof}

\begin{lemma}
	设 $\varepsilon_1,\cdots,\varepsilon_n$ 是 $(V,(\cdot,\cdot))$ 的一组标准正交基，且对任意 $\alpha,\beta\in V$，有
	\[
	\alpha=(\varepsilon_1,\cdots,\varepsilon_n)X,
	\qquad
	\beta=(\varepsilon_1,\cdots,\varepsilon_n)Y,
	\]
	其中
	\[
	X=\begin{pmatrix}
		x_1\\
		\vdots\\
		x_n
	\end{pmatrix},
	\qquad
	Y=\begin{pmatrix}
		y_1\\
		\vdots\\
		y_n
	\end{pmatrix}.
	\]
	则
	\[
	(\alpha,\beta)=X'Y.
	\]
\end{lemma}
\begin{proof}
	由坐标表示可知
	\[
	\alpha=\sum_{i=1}^{n}x_i\varepsilon_i,
	\qquad
	\beta=\sum_{j=1}^{n}y_j\varepsilon_j.
	\]
	于是
	\[
	\begin{aligned}
		(\alpha,\beta)
		&=\left(\sum_{i=1}^{n}x_i\varepsilon_i,
		\sum_{j=1}^{n}y_j\varepsilon_j\right)=\sum_{i,j}x_i y_j(\varepsilon_i,\varepsilon_j)=\sum_{i=1}^{n}x_i y_i=X'Y.
	\end{aligned}
	\]
\end{proof}

\section{度量映射与度量矩阵}



\begin{definition}
	设 $(V,(\cdot,\cdot)_V)$ 与 $(W,(\cdot,\cdot)_W)$ 是实内积空间，$f\in\Hom(V,W)$. 若对任意 $\alpha,\beta\in V$ 都有
	\[
	(f(\alpha),f(\beta))_W=(\alpha,\beta)_V,
	\]
	则称 $f$ 是一个\textbf{度量映射}. 特别地，取 $\beta=\alpha$ 得
	\[
	|f(\alpha)|=|\alpha|,
	\]
	故度量映射保持长度.
\end{definition}

\begin{proposition}
	若 $f:V\longrightarrow W$ 是度量映射，则 $f$ 必然是单射.
\end{proposition}
\begin{proof}
	任取 $\alpha\in\ker f$，则 $f(\alpha)=0$，从而
	\[
	(\alpha,\alpha)_V=(f(\alpha),f(\alpha))_W=(0,0)_W=0.
	\]
	由内积的正定性，$\alpha=0$，故 $\ker f=0$，于是 $f$ 是单射.
\end{proof}

\begin{definition}
	设 $V,W$ 是实内积空间. 若度量映射 $f:V\longrightarrow W$ 是满射，则称 $f$ 是从 $V$ 到 $W$ 的\textbf{等距}. 此时 $f$ 是线性同构且保持内积.
\end{definition}


度量映射不一定是满射. 例如
\[
\mathbb{R}^2\longrightarrow \mathbb{R}^3,
\qquad
\begin{pmatrix}x\\ y\end{pmatrix}
\longmapsto
\begin{pmatrix}x\\ y\\ 0\end{pmatrix}
\]
是度量映射，但不是满射.

\begin{definition}
	对内积空间 $V,W$，若存在从 $V$ 到 $W$ 的等距，则称 $V$ \textbf{等距于} $W$.
\end{definition}
	显然，等距关系是一个等价关系.


\begin{theorem}
	设 $V,W$ 是有限维欧氏空间，则 $V$ 与 $W$ 等距当且仅当$\dim V=\dim W.$
	
\end{theorem}
\begin{proof}
	若 $V$ 与 $W$ 等距，则存在从 $V$ 到 $W$ 的等距 $f$. 由上面的命题，$f$ 是线性同构，故 $\dim V=\dim W$.
	
	反之，设 $\dim V=\dim W=n$. 取 $V$ 的一组标准正交基 $\varepsilon_1,\cdots,\varepsilon_n$，取 $W$ 的一组标准正交基 $\eta_1,\cdots,\eta_n$. 由线性映射存在唯一性定理，存在线性映射
	\[
	f:V\longrightarrow W,
	\qquad
	f(\varepsilon_i)=\eta_i,
	\quad 1\leqslant i\leqslant n.
	\]
	则 $f$ 是线性同构. 对任意
	\[
	\alpha=(\varepsilon_1\cdots\varepsilon_n)X,
	\qquad
	\beta=(\varepsilon_1\cdots\varepsilon_n)Y,
	\]
	有
	\[
	f(\alpha)=(\eta_1\cdots\eta_n)X,
	\qquad
	f(\beta)=(\eta_1\cdots\eta_n)Y.
	\]
	由标准正交基下内积的坐标表达式，
	\[
	(\alpha,\beta)_V=X'Y=(f(\alpha),f(\beta))_W.
	\]
	故 $f$ 是等距.
\end{proof}



\begin{definition}
	设 $(V,(\cdot,\cdot))$ 是 $n$ 维欧氏空间，取 $V$ 的一组基 $\alpha_1,\cdots,\alpha_n$. 称矩阵
	\[
	A=((\alpha_i,\alpha_j))\in M_n(\mathbb{R})
	\]
	为内积 $(\cdot,\cdot)$ 在基 $\alpha_1,\cdots,\alpha_n$ 下的\textbf{度量矩阵}
\end{definition}

\begin{lemma}
	对上述矩阵 $A=((\alpha_i,\alpha_j))$，若
	\[
	\alpha=(\alpha_1\cdots\alpha_n)X,
	\qquad
	\beta=(\alpha_1\cdots\alpha_n)Y,
	\]
	其中 $X=(x_1,\cdots,x_n)'$，$Y=(y_1,\cdots,y_n)'$，则
	\[
	(\alpha,\beta)=X'AY.
	\]
\end{lemma}
\begin{proof}

	\[
	\begin{aligned}
		(\alpha,\beta)
		&=\left(\sum_{i=1}^n x_i\alpha_i,\sum_{j=1}^n y_j\alpha_j\right)             =\sum_{i,j}x_i y_j(\alpha_i,\alpha_j)         =\sum_{i,j}x_iA(i,j)y_j=X'AY.
	\end{aligned}
	\]
\end{proof}

\begin{proposition}
	欧氏空间的内积在任意基下的度量矩阵都是实对称可逆矩阵.
\end{proposition}
\begin{proof}
	设 $A=((\alpha_i,\alpha_j))$. 由内积的对称性，$A'=A$.
	
	若 $A$ 不可逆，则存在 $0\ne X_0\in\mathbb{R}^n$，使得 $AX_0=0$. 令
	\[
	\alpha=(\alpha_1\cdots\alpha_n)X_0.
	\]
	由于 $\alpha_1,\cdots,\alpha_n$ 是基且 $X_0\ne0$，故 $\alpha\ne0$. 但
	\[
	(\alpha,\alpha)=X_0'AX_0=0,
	\]
	这与正定性矛盾. 故 $A$ 可逆.
\end{proof}

\begin{proposition}
	设 $(V,(\cdot,\cdot))$ 是实内积空间，$\beta_1,\cdots,\beta_m\in V$. 记
	\[
	G(\beta_1,\cdots,\beta_m)=((\beta_i,\beta_j)).
	\]
	$G(\beta_1,\cdots,\beta_m)$称为$\beta_1,\cdots,\beta_m$的\textbf{Gram 矩阵}.
	
	则 $\beta_1,\cdots,\beta_m$ 线性无关当且仅当 $G(\beta_1,\cdots,\beta_m)$ 可逆.
\end{proposition}
\begin{proof}
	充分性：$G(\beta_1,\cdots,\beta_m)$ 可逆. 设
	\[
	x_1\beta_1+\cdots+x_m\beta_m=0.
	\]
	两边分别与 $\beta_1,\cdots,\beta_m$ 作内积，得到齐次线性方程组
	\[
	\left\{
	\begin{aligned}
		(\beta_1,\beta_1)x_1+\cdots+(\beta_m,\beta_1)x_m&=0,\\
		(\beta_1,\beta_2)x_1+\cdots+(\beta_m,\beta_2)x_m&=0,\\
		&\vdots\\
		(\beta_1,\beta_m)x_1+\cdots+(\beta_m,\beta_m)x_m&=0.
	\end{aligned}
	\right.
	\]
	该方程组仅有平凡解，故$\beta_1,\cdots,\beta_m$ 线性无关
	
	反之，若 $\beta_1,\cdots,\beta_m$ 线性无关，设 $W=\langle\beta_1,\cdots,\beta_m\rangle$，则 $W$ 是一个欧氏空间，且 $G(\beta_1,\cdots,\beta_m)$ 是内积 $(\cdot,\cdot)|_W$ 在基 $\beta_1,\cdots,\beta_m$ 下的度量矩阵。由前已证，欧氏空间的度量矩阵必然可逆，故 $G(\beta_1,\cdots,\beta_m)$ 可逆。
\end{proof}

\begin{example}
	设 $(V,(\cdot,\cdot))$ 是欧氏空间，$\alpha_1,\cdots,\alpha_n$ 和 $\beta_1,\cdots,\beta_n$ 是 $V$ 的两组基. 记
	\[
	A=((\alpha_i,\alpha_j)),
	\qquad
	B=((\beta_i,\beta_j)).
	\]
	若可逆矩阵 $C\in M_n(\mathbb{R})$ 满足
	\[
	(\beta_1\cdots\beta_n)=(\alpha_1\cdots\alpha_n)C,
	\]
	\[
	\begin{aligned}
		B(i,j)&=(\beta_i,\beta_j) =\left(\sum_{k=1}^n C(k,i)\alpha_k,\sum_{l=1}^n C(l,j)\alpha_l\right)=\sum_{k,l=1}^n C(k,i)C(l,j)(\alpha_k,\alpha_l)=(C'AC)(i,j).
	\end{aligned}
	\]
	因此 $B=C'AC$.
\end{example}


\begin{definition}
	设 $A,B\in M_n(\mathbb{F})$. 若存在 $C\in \GL_n(\mathbb{F})$，使得
	\[
	B=C'AC,
	\]
	则称 $A$ 与 $B$ \textbf{合同}，也称\textbf{相合}. 合同关系是一个等价关系.
\end{definition}

\begin{corollary}
	欧氏空间的内积在不同基下的度量矩阵合同.
\end{corollary}

\begin{example}
	设 $A$ 是欧氏空间的内积在某组基下的度量矩阵，则
	\begin{enumerate}
		\item $A$ 与单位阵 $E_n$ 合同；
		\item 对任意 $X\in\mathbb{R}^n$，都有 $X'AX\geqslant0$，且 $X'AX=0$ 当且仅当 $X=0$.
	\end{enumerate}
\end{example}
\begin{proof}
	取该欧氏空间的一组标准正交基，则内积在这组基下的度量矩阵为 $E_n$，由上面的推论知 $A$ 与 $E_n$ 合同.
	
	若 $\alpha=(\alpha_1\cdots\alpha_n)X$，则
	\[
	X'AX=(\alpha,\alpha)\geqslant0.
	\]
	并且 $X'AX=0$ 当且仅当 $(\alpha,\alpha)=0$，即 $\alpha=0$. 因 $\alpha_1,\cdots,\alpha_n$ 是基，这又等价于 $X=0$.
\end{proof}

\begin{definition}
	对实对称矩阵 $A\in M_n(\mathbb{R})$，若对任意 $X\in\mathbb{R}^n$ 都有
	\[
	X'AX\geqslant0,
	\]
	且
	\[
	X'AX=0\Longleftrightarrow X=0,
	\]
	则称 $A$ 是一个\textbf{正定阵}.
\end{definition}

\begin{example}
	设 $A=\operatorname{diag}(a_1,\cdots,a_n)$ 是实对角阵，则 $A$ 正定当且仅当
	\[
	a_i>0,
	\qquad 1\leqslant i\leqslant n.
	\]
	此时
	\[
	X'AX=a_1x_1^2+\cdots+a_nx_n^2.
	\]
\end{example}

\begin{example}
	设 $A,B$ 是 $n$ 阶实对称阵，且 $A,B$ 合同，则 $A$ 正定当且仅当 $B$ 正定.
\end{example}
\begin{proof}
	设
	\[
	B=C'AC,
	\qquad C\in \GL_n(\mathbb{R}).
	\]
	若 $A$ 正定，则对任意 $X\in\mathbb{R}^n$，有
	\[
	X'BX=X'C'ACX=(CX)'A(CX)\geqslant0.
	\]
	若 $X'BX=0$，则 $(CX)'A(CX)=0$，由 $A$ 正定得 $CX=0$，又 $C$ 可逆，故 $X=0$. 因此 $B$ 正定. 反向同理.
\end{proof}

\begin{proposition}
	设 $A\in S_n(\mathbb{R})$，则 $A$ 是正定阵当且仅当 $A$ 是某个 $n$ 维欧氏空间的内积在某组基下的度量矩阵.
\end{proposition}
\begin{proof}
	先设 $A$ 是正定阵. 取 $V$ 的一组基 $\alpha_1,\cdots,\alpha_n$. 对
	\[
	\alpha=(\alpha_1\cdots\alpha_n)X,
	\qquad
	\beta=(\alpha_1\cdots\alpha_n)Y,
	\]
	定义
	\[
	(\alpha,\beta)=X'AY.
	\]
	容易验证 $(\cdot,\cdot)$ 是 $V$ 上的一个内积，且 $A$ 正是该内积在基 $\alpha_1,\cdots,\alpha_n$ 下的度量矩阵.
	
	反之，若 $A$ 是某个欧氏空间的内积在某组基下的度量矩阵，则对任意 $X\in\mathbb{R}^n$，$X\neq0$，令 $\alpha=(\alpha_1\cdots\alpha_n)X\neq0$，有 $X'AX=(\alpha,\alpha)>0$，故 $A$ 正定。
\end{proof}

\begin{corollary}
	设 $A$ 是 $n$ 阶实对称阵，则下列条件等价：
	\begin{enumerate}
		\item $A$ 正定；
		\item $A$ 与单位阵 $E_n$ 合同；
		\item 存在实可逆矩阵 $C$，使得
		\[
		A=C'C.
		\]
	\end{enumerate}
\end{corollary}


\begin{example}
	设 $V$ 是 $C[a,b]$ 的一个有限维子空间，$\varphi_1(x),\cdots,\varphi_n(x)$ 是 $V$ 的一组基. 设矩阵 $A$ 的 $(i,j)$ 元为
	\[
	A(i,j)=\int_a^b \varphi_i(x)\varphi_j(x)\,dx,
	\qquad 1\leqslant i,j\leqslant n.
	\]
	则 $A$ 是正定阵.
\end{example}
\begin{example}
	设矩阵 $A$ 的 $(i,j)$ 元为
	\[
	A(i,j)=\frac{1}{i+j+1},
	\qquad 1\leqslant i,j\leqslant n.
	\]
	则 $A$ 是正定阵.
\end{example}

\section{欧氏空间上的正交变换与正交矩阵}

\begin{definition}
	设 $(V,(\cdot,\cdot))$ 是 $n$ 维欧氏空间. 若 $T\in\End(V)$ 满足
	\[
	(T(\alpha),T(\beta))=(\alpha,\beta),\qquad \forall\alpha,\beta\in V,
	\]
	则称 $T$ 是 $V$ 上的一个\textbf{正交变换}.
\end{definition}

\begin{example}
	设 $(V,(\cdot,\cdot))$ 是欧氏空间，则恒等映射 $id_V$ 是 $V$ 的正交变换.
\end{example}

\begin{example}
	设 $(V,(\cdot,\cdot))$ 是欧氏空间，记
	\[
	O(V)=\{T\in \GL(V):T\text{ 是 }V\text{ 上的正交变换}\}.
	\]
	则 $O(V)$ 是 $\GL(V)$ 的子群，称为 $V$ 的\textbf{正交群}.
\end{example}

\begin{definition}
	设 $(V,(\cdot,\cdot))$ 是欧氏空间，任意取定 $0\ne\alpha\in V$.定义
	\[
	r_{\alpha}:V\longrightarrow V,
	\qquad
	r_{\alpha}(\beta)=\beta-2\frac{(\beta,\alpha)}{(\alpha,\alpha)}\alpha,
	\]
	 称 $r_{\alpha}$ 为由 $\alpha$ 所给出的\textbf{反射变换}.
\end{definition}
\begin{proposition}
	上述$r_{\alpha}\in O(V)$.
\end{proposition}
\begin{proof}
	首先 $r_{\alpha}$ 是线性映射，因为
	\[
	\begin{aligned}
		r_{\alpha}(k_1\beta_1+k_2\beta_2)
		&=k_1\beta_1+k_2\beta_2
		-2\frac{(k_1\beta_1+k_2\beta_2,\alpha)}{(\alpha,\alpha)}\alpha \\
		&=k_1r_{\alpha}(\beta_1)+k_2r_{\alpha}(\beta_2).
	\end{aligned}
	\]
	又对任意 $\beta_1,\beta_2\in V$，有
	\[
	\begin{aligned}
		(r_\alpha(\beta_1),r_\alpha(\beta_2))
		&=\left(\beta_1-2\frac{(\beta_1,\alpha)}{(\alpha,\alpha)}\alpha,
		\beta_2-2\frac{(\beta_2,\alpha)}{(\alpha,\alpha)}\alpha\right)\\
		&=(\beta_1,\beta_2)-2\frac{(\beta_1,\alpha)(\beta_2,\alpha)}{(\alpha,\alpha)}
		-2\frac{(\beta_1,\alpha)(\beta_2,\alpha)}{(\alpha,\alpha)}
		+4\frac{(\beta_1,\alpha)(\beta_2,\alpha)}{(\alpha,\alpha)}\\
		&=(\beta_1,\beta_2).
	\end{aligned}
	\]
	所以 $r_\alpha\in O(V)$.
\end{proof}

\begin{proposition}
	设 $0\ne\alpha\in V$，则反射变换 $r_\alpha$ 满足
	\begin{enumerate}
		\item $r_\alpha^2=id_V$；
		\item $r_\alpha$ 可对角化；
	\end{enumerate}
\end{proposition}
\begin{proof}
	设
	\[r_\alpha(\beta)
	=
	\beta-2\frac{(\beta,\alpha)}{(\alpha,\alpha)}\alpha,
	\qquad \beta\in V.\]

	由正交分解$ V=\langle \alpha\rangle\oplus \langle \alpha\rangle^\perp,$
	任意 $\beta\in V$ 可唯一写为$\beta=k\alpha+\gamma,\gamma\in \langle \alpha\rangle^\perp. $于是
	\[	r_\alpha(\beta)
	=
	k\alpha+\gamma-2\frac{(k\alpha+\gamma,\alpha)}{(\alpha,\alpha)}\alpha
	=
	-k\alpha+\gamma.\]
再作用一次得
\[r_\alpha^2(\beta)=r_\alpha(-k\alpha+\gamma)=k\alpha+\gamma=\beta.\]
	故 $r_\alpha^2=id_V$.
	
同时，由上式可知
	\[
	r_\alpha(\alpha)=-\alpha,
	\qquad
	r_\alpha(\gamma)=\gamma,\quad \gamma\in \langle \alpha\rangle^\perp.
	\]
	因此 $\langle \alpha\rangle$ 是 $r_\alpha$ 关于特征值 $-1$ 的特征子空间，
	$\langle \alpha\rangle^\perp$ 是 $r_\alpha$ 关于特征值 $1$ 的特征子空间.
	又
	\[
	V=\langle \alpha\rangle\oplus \langle \alpha\rangle^\perp,
	\]
	所以 $r_\alpha$ 可对角化.
	
	
	
\end{proof}


\begin{lemma}
	设 $(V,(\cdot,\cdot))$ 是欧氏空间，$T\in O(V)$. 若 $0\ne W$ 是 $T$-子空间，则
	\[
	T|_W\in O(W).
	\]
\end{lemma}
\begin{proof}
	由 $W$ 是 $T$-子空间，$T(W)\subset W$. 又 $T$ 可逆，且 $V$ 有限维，故 $T|_W$ 是 $W$ 上的线性同构. 对任意 $\alpha,
	\beta\in W$，有
	\[
	(T|_W(\alpha),T|_W(\beta))=(T\alpha,T\beta)=(\alpha,\beta),
	\]
	所以 $T|_W\in O(W)$.
\end{proof}

\begin{corollary}\label{T-子空间}
	设 $(V,(\cdot,\cdot))$ 是欧氏空间，$T\in O(V)$. 若 $W$ 是 $T$-子空间，则 $W^\perp$ 也是 $T$-子空间.
\end{corollary}
\begin{proof}
	任取 $\alpha\in W^\perp$ 与 $w\in W$. 由上一引理，$T|_W\in O(W)$，故 $T(W)=W$ 且 $T^{-1}w\in W$. 因而
	\[
	(T\alpha,w)=(T\alpha,T(T^{-1}w))=(\alpha,T^{-1}w)=0.
	\]
	所以 $T\alpha\in W^\perp$.
\end{proof}

\begin{proposition}
	(Cartan--Dieudonn\'e 定理)
	任意正交变换都是若干个反射变换的乘积.
\end{proposition}
\begin{proof}
	 对 $\dim V$ 作归纳.
	当 $\dim V=1$ 时，任意正交变换只能为 $id_V$ 或某个反射变换，结论成立.
	
	设 $\dim V=n$ 时结论成立，考虑 $\dim V=n+1$ 的情形.
	若 $T=id_V$，则 $T$ 是零个反射变换的乘积. 若 $T\neq id_V$，取 $x\in V$ 使得
	$Tx\neq x$，令$u=x-Tx\neq 0. $因为 $T$ 是正交变换，所以$(x,x)=(Tx,Tx).$
	于是
	\[
	(u,u)=(x-Tx,x-Tx)=2(x,x)-2(x,Tx),
	\]
	且$(x,u)=(x,x)-(x,Tx). $
	因此
	\[
	r_u(x)=x-2\frac{(x,u)}{(u,u)}ux-u=
	Tx.\]
	令$S=r_uT. $则
	\[
	Sx=r_u(Tx)=x,
	\]
	即 $S$ 固定一维子空间 $\langle x\rangle$.
	又 $S$ 是正交变换，所以,由\ref{T-子空间}$, W=\langle x\rangle^\perp$
	是 $S$ 的不变子空间，并且 $S|_W\in O(W)$.
	
	由归纳假设，$S|_W$ 是 $W$ 上若干个反射变换的乘积. 这些反射变换都可以看成
	$V$ 上的反射变换，并且在 $\langle x\rangle$ 上恒等. 因此 $S$ 是 $V$ 上若干个反射
	变换的乘积. 又因为
	\[
	T=r_uS,
	\]
	故 $T$ 也是若干个反射变换的乘积.
	
\end{proof}
\begin{example}
	设 $(V,(\cdot,\cdot))$ 是欧氏空间，$T\in O(V)$. 设 $\varepsilon_1,\cdots,\varepsilon_n$ 是 $V$ 的一组标准正交基，并设
	\[
	(T(\varepsilon_1)\ \cdots\ T(\varepsilon_n))=(\varepsilon_1\ \cdots\ \varepsilon_n)A.
	\]
	则
	\[
	A'A=E_n.
	\]
\end{example}
\begin{proof}
	记 $A=(A(i,j))$. 由于
	\[
	T(\varepsilon_j)=\sum_{k=1}^{n}A(k,j)\varepsilon_k,
	\]
	故
	\[
	\begin{aligned}
		\delta_{ij}
		&=(\varepsilon_i,\varepsilon_j)=(T(\varepsilon_i),T(\varepsilon_j))=\left(\sum_{k=1}^{n}A(k,i)\varepsilon_k,
		\sum_{\ell=1}^{n}A(\ell,j)\varepsilon_\ell\right)\\
		&=\sum_{k,\ell=1}^{n}A(k,i)A(\ell,j)(\varepsilon_k,\varepsilon_\ell)
		=\sum_{k=1}^{n}A(k,i)A(k,j)=(A'A)(i,j).
	\end{aligned}
	\]
	所以 $A'A=E_n$.
\end{proof}

\begin{definition}
	设 $B\in M_n(\mathbb{R})$. 若
	\[
	B'B=E_n,
	\]
	则称 $B$ 是\textbf{正交矩阵}.
\end{definition}

由定义可知，若 $B$ 是正交矩阵，则
\[
B'B=E_n\Longleftrightarrow BB'=E_n\Longleftrightarrow B^{-1}=B'.
\]
准对角阵$\operatorname{diag}(A_1,\cdots,A_s)$是正交阵当且仅当每一个$A_i$是正交阵
\begin{lemma}
	设 $(V,(\cdot,\cdot))$ 是 $n$ 维欧氏空间，$T\in\End(V)$. 则 $T\in O(V)$ 当且仅当 $T$ 在任意一组标准正交基下的矩阵都是正交矩阵.
\end{lemma}
\begin{proof}
	必要性由上例已证. 反之，设 $T$ 在某一组标准正交基 $\varepsilon_1,
	\cdots,\varepsilon_n$ 下的矩阵为 $A$，且 $A'A=E_n$. 对任意
	\[
	\alpha=(\varepsilon_1\cdots\varepsilon_n)X,
	\qquad
	\beta=(\varepsilon_1\cdots\varepsilon_n)Y,
	\]
	有
	\[
	(T\alpha,T\beta)=(AX)'(AY)=X'A'AY=X'Y=(\alpha,\beta).
	\]
	故 $T\in O(V)$.
\end{proof}

\begin{corollary}
	设 $(V,(\cdot,\cdot))$ 是欧氏空间，$T\in\End(V)$. 则以下四条等价：
	\begin{enumerate}
		\item $T\in O(V)$；
		\item 对任意 $\alpha\in V$ 都有 $|T\alpha|=|\alpha|$；
		\item $T$ 把标准正交基变为标准正交基；
		\item $T$ 在任意标准正交基下的矩阵是正交矩阵.
	\end{enumerate}
\end{corollary}
\begin{proof}
	$(1)\Rightarrow(2)$ 由 $(T\alpha,T\alpha)=(\alpha,\alpha)$ 即得.
	$(2)\Rightarrow(1)$：对任意 $\alpha,\beta\in V$，利用
	\[
	(T(\alpha+\beta),T(\alpha+\beta))=(\alpha+\beta,\alpha+\beta)
	\]
	并展开，结合 $|T\alpha|=|\alpha|, |T\beta|=|\beta|$，得
	\[
	(T\alpha,T\beta)=(\alpha,\beta).
	\]
	$(1)\Leftrightarrow(4)$ 由上一引理得到. $(1)\Rightarrow(3)$ 显然；$(3)\Rightarrow(1)$ 时，若
	\[
	\alpha=(\varepsilon_1\cdots\varepsilon_n)X,
	\qquad
	\beta=(\varepsilon_1\cdots\varepsilon_n)Y,
	\]
	则
	\[
	(T\alpha,T\beta)=X'Y=(\alpha,\beta).
	\]
\end{proof}

\begin{definition}
	记
	\[
	O_n(\mathbb{R})=\{A\in M_n(\mathbb{R}):A'A=E_n\},
	\]
	称为 $\mathbb{R}$ 上的 $n$ 阶\textbf{正交群}.
\end{definition}
	$O_n(\mathbb{R})\text{是} \GL_n(\mathbb{R})$的子群


\begin{proposition}
		设 $(V,(\cdot,\cdot))$ 是 $n$ 维欧氏空间，则有群同构
	\[
	O(V)\cong O_n(\mathbb{R}).
	\]
\end{proposition}

\begin{proof}
	任意取定 $V$ 的一组标准正交基 $\varepsilon_1,
	\cdots,\varepsilon_n$. 定义映射
	\[
	\rho:O(V)\longrightarrow O_n(\mathbb{R}),
	\qquad
	T\longmapsto T\text{在基}\varepsilon_1,
		\cdots,\varepsilon_n\text{的矩阵}A
	\]
	显然$\rho$是一个群同构
\end{proof}

\begin{proposition}
	设 $A\in M_n(\mathbb{R})$，$(V,(\cdot,\cdot))$ 是 $n$ 维欧氏空间. 以下六条等价：
	\begin{enumerate}
		\item $A\in O_n(\mathbb{R})$；
		\item $A$ 是 $\mathbb{R}^n$ 在标准内积下的正交变换 $X\mapsto AX$；
		\item $A$ 的列向量是 $\mathbb{R}^n$ 在标准内积下的一组标准正交基；
		\item $A$ 的行向量是 $\mathbb{R}^n$ 在标准内积下的一组标准正交基；
		\item $A$ 是 $V$ 的两组标准正交基之间的转移矩阵；
		\item $A$ 是 $V$ 上某个正交变换在某组标准正交基下的矩阵.
	\end{enumerate}
\end{proposition}
\begin{proof}
	$(1)\Leftrightarrow(2)$：任取 $X,Y\in\mathbb{R}^n$，有
	\[
	(AX,AY)=(AX)'(AY)=X'(A'A)Y.
	\]
	故保持标准内积当且仅当 $A'A=E_n$.
	
	$(1)\Leftrightarrow(3)$：设 $A=(\alpha_1\ \alpha_2\ \cdots\ \alpha_n)$，则
	\[
	A'A=\bigl((\alpha_i,\alpha_j)\bigr)_{n\times n}.
	\]
	所以 $A'A=E_n$ 当且仅当 $\alpha_1,
	\cdots,\alpha_n$ 为标准正交基.
	
	$(1)\Leftrightarrow(4)$：将 $A$ 按行写为
	\[
	A=\begin{pmatrix}
		\beta_1\\
		\beta_2\\
		\vdots\\
		\beta_n
	\end{pmatrix},
	\]
	则
	\[
	AA'=\bigl((\beta_i,\beta_j)\bigr)_{n\times n}.
	\]
	又 $A'A=E_n\Longleftrightarrow AA'=E_n$，故等价.
	
	$(1)\Rightarrow(5)$：取 $V$ 的一组标准正交基 $\varepsilon_1,
	\cdots,\varepsilon_n$，令
	\[
	(\eta_1\cdots\eta_n)=(\varepsilon_1\cdots\varepsilon_n)A.
	\]
	则
	\[
	\bigl((\eta_i,\eta_j)\bigr)=A'E_nA=A'A=E_n,
	\]
	故 $\eta_1,
	\cdots,\eta_n$ 是标准正交基.
	
	$(5)\Rightarrow(1)$：若 $\varepsilon_1,
	\cdots,\varepsilon_n$ 与 $\eta_1,
	\cdots,\eta_n$ 都是标准正交基，且
	\[
	(\eta_1\cdots\eta_n)=(\varepsilon_1\cdots\varepsilon_n)A,
	\]
	则
	\[
	E_n=\bigl((\eta_i,\eta_j)\bigr)=A'E_nA=A'A.
	\]
	$(1)\Leftrightarrow(6)$ 由正交变换与正交矩阵的关系即得.
\end{proof}

\begin{lemma}
	设 $M_1,M_2\in M_n(\mathbb{F})$. 则
	\[
	M_1=M_2
	\]
	当且仅当对任意 $X,Y\in\mathbb{F}^n$ 都有
	\[
	X'M_1Y=X'M_2Y.
	\]
\end{lemma}
\begin{proof}
	只需证充分性. 令 $M=M_1-M_2$，则对任意 $X,Y$ 都有 $X'MY=0$. 固定任意 $X_0\in\mathbb{F}^n$，于是
	\[
	X'MX_0=0,
	\qquad \forall X\in\mathbb{F}^n.
	\]
	这说明 $MX_0=0$. 由于 $X_0$ 任意，故 $M=0$，即 $M_1=M_2$.
\end{proof}

\begin{definition}
	记
	\[
	SO_n(\mathbb{R})=\{A\in O_n(\mathbb{R}):\det A=1\},
	\]
	称为 $n$ 阶\textbf{特殊正交群}.
\end{definition}

\begin{lemma}
	设 $A\in O_n(\mathbb{R})$，则
	\begin{enumerate}
		\item $\det A=\pm1$，且 $A^{-1}=A'$；
		\item $A$ 的特征值都位于复平面上的单位圆上.
	\end{enumerate}
\end{lemma}
\begin{proof}
	\begin{enumerate}
		\item 由 $A'A=E_n$，两边取行列式得 $(\det A)^2=1$，故 $\det A=\pm1$，且 $A^{-1}=A'$.
		\item 设 $A\xi=\lambda\xi$，其中 $0\ne\xi\in\mathbb{C}^n$.取共轭$\overline{A} \overline{\xi}= \overline{\lambda}\overline{\xi}$,由$A\in O_n(\R)$,$A\overline{\xi}=\overline{\lambda}\overline{\xi}$,从而同时乘$\xi'$
		\[\xi'(A\bar{\xi}) = \bar{\lambda}\xi'\bar{\xi}\rightarrow (\xi'A)\bar{\xi} = \bar{\lambda}\xi'\bar{\xi}\]
		而$\xi'A = (A'\xi)'$,且
		\[\lambda \neq 0,\xi = A'A\xi = \lambda A'\xi \text{从而}A'\xi = \frac{1}{\lambda}\xi\]
		故\[\xi'\bar{\xi} =\bar{\lambda}\lambda\bar{\xi}\xi'\]
		$\lambda\bar{\lambda} =1 \rightarrow |\lambda|=1$
	\end{enumerate}
	

\end{proof}

\begin{lemma}
	设 $V$ 是任意的有限维实线性空间，$T\in\End(V)$. 则 $T$ 有一维或二维的不变子空间.
\end{lemma}
\begin{proof}
	若 $T$ 有实特征值，则其特征向量张成的一维子空间即为不变子空间.
	若 $T$ 没有实特征值，则取特征多项式的一个非实复根 $a+bi$，其中 $b>0$. 设对应复特征向量为 $\xi+i\eta$，其中 $\xi,\eta\in V$，不全为零. 由
	\[
	T(\xi+i\eta)=(a+bi)(\xi+i\eta)
	\]
	比较实部与虚部，得
	\[
	T\xi=a\xi-b\eta,
	\qquad
	T\eta=b\xi+a\eta.
	\]
	故 $\langle\xi,\eta\rangle$ 是 $T$ 的二维不变子空间.
\end{proof}

\begin{lemma}
	设 $(V,(\cdot,\cdot))$ 是欧氏空间，$T\in O(V)$. 则
	\[
	V=U_1\oplus\cdots\oplus U_r\oplus W_1\oplus\cdots\oplus W_s,
	\]
	其中这些子空间两两正交，且 $U_i$ 是一维 $T$-子空间，$W_j$ 是二维 $T$-子空间，并且 $W_j$ 不能再分解为一维 $T$-子空间的直和.
\end{lemma}
\begin{proof}
	对 $\dim V$ 作归纳. 当 $\dim V=1$ 时显然. 若 $\dim V>1$，由上一引理，取 $T$ 的一个一维或二维不变子空间 $W$. 因 $T\in O(V)$，由前面的推论，$W^\perp$ 也是 $T$-子空间，且
	\[
	V=W\oplus W^\perp.
	\]
	对 $W^\perp$ 使用归纳假设即可.
\end{proof}

\begin{corollary}
	设 $(V,(\cdot,\cdot))$ 是欧氏空间，$T\in O(V)$. 则存在 $V$ 的一组标准正交基，使得 $T$ 在这组标准正交基下的矩阵是由若干一阶正交阵与二阶正交阵构成的准对角阵；其中，出现的二阶正交阵没有实特征值.
\end{corollary}
\begin{corollary}
	设$A\in O_n(\R)$,存在正交矩阵 $U$使得$ U^{-1}AU$是准对角阵，由若干一阶正交阵与二阶正交阵构成的准对角阵；其中，出现的二阶正交阵没有实特征值.

\end{corollary}
\begin{definition}
	设 $A,B\in M_n(\mathbb{R})$. 若存在正交矩阵 $U$ 使得
	\[
	U^{-1}AU=B,
	\]
	则称 $A$ 与 $B$ \textbf{正交相似}.
\end{definition}
\begin{lemma}
	$A$ 是实对称矩阵，不同特征值对应的特征向量正交
\end{lemma}
\begin{proof}
	设 $A'=A$，且
	\[
	A\alpha=\lambda\alpha,\qquad A\beta=\mu\beta,\qquad \lambda\ne\mu.
	\]
	由于 $A'=A$，故
	\[
	(A\alpha,\beta)=(\alpha,A\beta).
	\]
	于是
	\[
	\lambda(\alpha,\beta)=\mu(\alpha,\beta).
	\]
	从而
	\[
	(\lambda-\mu)(\alpha,\beta)=0.
	\]
	因为 $\lambda\ne\mu$，所以
	\[
	(\alpha,\beta)=0.
	\]
	故不同特征值对应的特征向量正交.
\end{proof}
\begin{lemma}
	任意二阶正交矩阵 $A$ 都正交相似于如下四个两两不相似的正交矩阵之一：
		\[
		\begin{pmatrix}
			1&0\\
			0&1
		\end{pmatrix},
		\quad
		\begin{pmatrix}
			-1&0\\
			0&-1
		\end{pmatrix},
		\quad
		\begin{pmatrix}
			1&0\\
			0&-1
		\end{pmatrix},
		\quad
		\begin{pmatrix}
			\cos\theta&\sin\theta\\
			-\sin\theta&\cos\theta
		\end{pmatrix},
		\quad 0<\theta<\pi .
		\]
\end{lemma}

\begin{proof}
	设
	\[
	A=
	\begin{pmatrix}
		a&b\\
		c&d
	\end{pmatrix}
	\in O_2(\mathbb R).
	\]
	由 $AA'=E_2$ 得
	\[
	a^2+b^2=1,\qquad ac+bd=0,\qquad c^2+d^2=1.
	\]
	因此 $(a,b),(c,d)$ 都在单位圆上，且相互垂直.
	
	设
	\[
	(a,b)=(\cos\theta,\sin\theta),
	\qquad 0\leq \theta<2\pi.
	\]
	由于 $(c,d)$ 是与 $(a,b)$ 垂直的单位向量，故
	\[
	(c,d)=(-\sin\theta,\cos\theta)
	\]
	或
	\[
	(c,d)=(\sin\theta,-\cos\theta).
	\]
	于是
	\[
	A=
	\begin{pmatrix}
		\cos\theta&\sin\theta\\
		-\sin\theta&\cos\theta
	\end{pmatrix}
	\]
	或
	\[
	A=
	\begin{pmatrix}
		\cos\theta&\sin\theta\\
		\sin\theta&-\cos\theta
	\end{pmatrix}.
	\]
	
	先考虑第二种情形. 
	显然 $A'=A$，即 $A$ 是实对称阵. 又因为 $A\in O_2(\mathbb R)$，所以$	A'A=E_2.$由 $A'=A$ 得$A^2=E_2.$因此，若 $\lambda$ 是 $A$ 的特征值，取对应的非零特征向量 $\xi$，则$	A\xi=\lambda\xi.$两边再作用一次 $A$，得到$	A^2\xi=\lambda^2\xi.$而 $A^2=E_2$，故$\xi=\lambda^2\xi.$由于 $\xi\ne 0$，所以$\lambda^2=1.$于是$\lambda=\pm 1.$
	
	
	另一方面，
	\[
	\det A
	=
	-\cos^2\theta-\sin^2\theta
	=
	-1.
	\]
	所以 $A$ 的两个特征值分别为 $1$ 和 $-1$.
	
	设 $\xi_1,\xi_2$ 分别是对应于特征值 $1,-1$ 的单位特征向量. 
	因为 $A$ 是实对称阵，所以不同特征值对应的特征向量正交. 
	于是 $\xi_1,\xi_2$ 构成 $\mathbb R^2$ 的一个标准正交基.
	
	在这组标准正交基下，$A$ 的矩阵为
	\[
	\begin{pmatrix}
		1&0\\
		0&-1
	\end{pmatrix}.
	\]
	故 $A$ 正交相似于
	\[
	\begin{pmatrix}
		1&0\\
		0&-1
	\end{pmatrix}.
	\]
	
	再考虑第一种情形，记
	\[
	R_\theta=
	\begin{pmatrix}
		\cos\theta&\sin\theta\\
		-\sin\theta&\cos\theta
	\end{pmatrix}.
	\]
	若 $\theta=0$，则
	\[
	R_\theta=E_2.
	\]
	若 $\theta=\pi$，则
	\[
	R_\theta=-E_2.
	\]
	若 $0<\theta<\pi$，则得到第四类矩阵.
	
	若 $\pi<\theta<2\pi$，令
	\[
	\theta'=2\pi-\theta,
	\qquad 0<\theta'<\pi.
	\]
	则
	\[
	R_\theta=
	\begin{pmatrix}
		\cos\theta'&-\sin\theta'\\
		\sin\theta'&\cos\theta'
	\end{pmatrix}.
	\]
	取
	\[
	P=
	\begin{pmatrix}
		0&1\\
		1&0
	\end{pmatrix},
	\qquad P'=P,\quad P^2=E_2.
	\]
	则 $P\in O_2(\mathbb R)$，且
	\[
	P
	\begin{pmatrix}
		\cos\theta'&-\sin\theta'\\
		\sin\theta'&\cos\theta'
	\end{pmatrix}
	P
	=
	\begin{pmatrix}
		\cos\theta'&\sin\theta'\\
		-\sin\theta'&\cos\theta'
	\end{pmatrix}.
	\]
	故当 $\pi<\theta<2\pi$ 时，$R_\theta$ 也正交相似于某个
	\[
	\begin{pmatrix}
		\cos\theta'&\sin\theta'\\
		-\sin\theta'&\cos\theta'
	\end{pmatrix},
	\qquad 0<\theta'<\pi.
	\]
	
	最后说明这些标准形式两两不相似.
	前三个矩阵的特征值分别为
	\[
	1,1;\qquad -1,-1;\qquad 1,-1,
	\]
	故两两不相似. 而
	\[
	\begin{pmatrix}
		\cos\theta&\sin\theta\\
		-\sin\theta&\cos\theta
	\end{pmatrix}
	\]
	的特征多项式为
	\[
	\lambda^2-2\cos\theta\,\lambda+1.
	\]
	当 $0<\theta<\pi$ 时，其特征值为
	\[
	\cos\theta\pm i\sin\theta,
	\]
	不是实数，故它不与前三类矩阵相似.
	
	若 $0<\theta_1,\theta_2<\pi$ 且 $\theta_1\ne \theta_2$，由于 $\cos x$ 在 $(0,\pi)$ 上单调，
	所以
	\[
	\cos\theta_1\ne \cos\theta_2.
	\]
	从而对应的特征多项式不同，故
	\[
	R_{\theta_1}\nsim R_{\theta_2}.
	\]
	综上，任意二阶正交矩阵都正交相似于上述四类之一，且这些标准形式两两不相似.
\end{proof}

\begin{theorem}
	设 $A\in M_n(\mathbb{R})$. 则 $A\in O_n(\mathbb{R})$ 当且仅当存在正交矩阵 $U$，使得
	\[
	U^{-1}AU=	\operatorname{diag}\left((a_1),\cdots,(a_r),
	\begin{pmatrix}\cos\theta_1&\sin\theta_1\\-\sin\theta_1&\cos\theta_1\end{pmatrix},
	\cdots,
	\begin{pmatrix}\cos\theta_s&\sin\theta_s\\-\sin\theta_s&\cos\theta_s\end{pmatrix}
	\right),
	\]
	其中
	\[
	a_i=1\text{ 或 }-1,
	\qquad
	0<\theta_i<\pi,
	\qquad
	r+2s=n.
	\]
\end{theorem}

\begin{corollary}
	任意正交矩阵在 $\mathbb{R}$ 上的初等因子只能形如
	\[
	\lambda+1,
	\qquad
	\lambda-1,
	\qquad
	\lambda^2-2\lambda\cos\theta+1,
	\quad 0<\theta<\pi.
	\]
\end{corollary}

\begin{example}
	设 $A\in O_n(\mathbb{R})$，则
	\[
	|\operatorname{tr}(A)|\leqslant n.
	\]
	进一步地，若 $\operatorname{tr}(A)=n$，则 $A=E_n$；若 $\operatorname{tr}(A)=-n$，则 $A=-E_n$.
\end{example}
\begin{proof}
	由上一定理，$A$ 正交相似于由 $\pm1$ 与旋转块构成的准对角矩阵. 旋转块
	\[
	\begin{pmatrix}\cos\theta&\sin\theta\\-\sin\theta&\cos\theta\end{pmatrix}
	\]
	对迹的贡献为 $2\cos\theta$，其绝对值不超过 $2$. 因此 $|\operatorname{tr}(A)|\leq n$. 若 $\operatorname{tr}(A)=n$，则每个一阶块均为 $1$，每个二阶块均须满足 $2\cos\theta=2$，这与 $0<\theta<\pi$ 矛盾，故没有二阶块且 $A=E_n$. $\operatorname{tr}(A)=-n$ 同理.
\end{proof}

\begin{example}
	设
	\[
	A=\begin{pmatrix}
		\dfrac23&\dfrac23&\dfrac13\\[4pt]
		\dfrac13&-\dfrac23&\dfrac23\\[4pt]
		-\dfrac23&\dfrac13&\dfrac23
	\end{pmatrix}.
	\]
	求正交矩阵 $U$，使得 $U^{-1}AU$ 由一个一阶方阵和一个二阶方阵构成的准对角阵.
\end{example}
\begin{proof}[解答]
	直接展开得
	\[
	f(\lambda)=|\lambda E_3-A|
	=\lambda^3-\frac23\lambda^2-\frac23\lambda+1
	=(\lambda+1)\left(\lambda^2-\frac53\lambda+1\right).
	\]
	因此一阶块为 $(-1)$，二阶旋转块满足
	\[
	2\cos\theta=\frac53,
	\qquad
	\cos\theta=\frac56,
	\qquad
	\sin\theta=\frac{\sqrt{11}}6.
	\]
	所以存在正交矩阵 $U$ 使得
	\[
	U^{-1}AU=\begin{pmatrix}
		-1&0&0\\
		0&\dfrac56&\dfrac{\sqrt{11}}6\\[4pt]
		0&-\dfrac{\sqrt{11}}6&\dfrac56
	\end{pmatrix}.
	\]
\end{proof}


\section{对称变换与实对称阵}

\begin{definition}
	设 $(V,(\cdot,\cdot))$ 是实内积空间，$T\in\End(V)$. 如果对任意 $\alpha,\beta\in V$ 都有
	\[
	(T(\alpha),\beta)=(\alpha,T(\beta)),
	\]
	则称 $T$ 是 $V$ 的一个\textbf{对称变换}(symmetric transformation).
\end{definition}

记号：用 $\mathcal{S}(V)$ 表示 $V$ 的所有对称变换所组成的集合. 显然，对任意 $a\in\mathbb{R}$，有
\[
a\id_V\in\mathcal{S}(V).
\]

\begin{lemma}
	设 $(V,(\cdot,\cdot))$ 是实内积空间. 则 $\mathcal{S}(V)$ 是 $\End(V)$ 的子空间.
\end{lemma}
\begin{proof}
	设 $T_1,T_2\in\mathcal{S}(V)$，$a_1,a_2\in\mathbb{R}$. 对任意 $\alpha,\beta\in V$，有
	\[
	\begin{aligned}
		((a_1T_1+a_2T_2)(\alpha),\beta)
		&=(a_1T_1(\alpha),\beta)+(a_2T_2(\alpha),\beta)\\
		&=a_1(\alpha,T_1(\beta))+a_2(\alpha,T_2(\beta))\\
		&=(\alpha,(a_1T_1+a_2T_2)(\beta)).
	\end{aligned}
	\]
	故 $a_1T_1+a_2T_2\in\mathcal{S}(V)$.
\end{proof}

\begin{lemma}
	设 $(V,(\cdot,\cdot))$ 是实内积空间，$T\in\mathcal{S}(V)$. 如果 $W$ 是 $T$-子空间，则 $W^\perp$ 也是 $T$-子空间.
\end{lemma}
\begin{proof}
	任取 $\alpha\in W^\perp$，要证 $T(\alpha)\in W^\perp$. 对任意 $\beta\in W$，由于 $W$ 是 $T$-子空间，有 $T(\beta)\in W$，于是
	\[
	(\beta,T(\alpha))=(T(\beta),\alpha)=0.
	\]
	故 $T(\alpha)\in W^\perp$.
\end{proof}

\begin{example}
	设 $(V,(\cdot,\cdot))$ 是实内积空间，$T\in\mathcal{S}(V)$. 设 $W$ 是非零的 $T$-子空间，则
	\[
	T|_W\in\mathcal{S}(W).
	\]
\end{example}
\begin{proof}
	对任意 $\alpha,\beta\in W$，有
	\[
	(T|_W(\alpha),\beta)=(T(\alpha),\beta)=(\alpha,T(\beta))=(\alpha,T|_W(\beta)).
	\]
	故 $T|_W\in\mathcal{S}(W)$.
\end{proof}

\begin{lemma}
	设 $(V,(\cdot,\cdot))$ 是实内积空间，$T\in\mathcal{S}(V)$. 设 $\lambda_1,\lambda_2$ 是 $T$ 的不同的特征值. 则属于不同特征值的特征向量彼此正交. 即若
	\[
	T(\xi_1)=\lambda_1\xi_1,
	\qquad
	T(\xi_2)=\lambda_2\xi_2,
	\qquad
	\lambda_1\ne\lambda_2,
	\]
	则 $(\xi_1,\xi_2)=0$.
\end{lemma}
\begin{proof}
	由 $T$ 是对称变换，有
	\[
	\lambda_1(\xi_1,\xi_2)=(\lambda_1\xi_1,\xi_2)=(T(\xi_1),\xi_2)=(\xi_1,T(\xi_2))=(\xi_1,\lambda_2\xi_2)=\lambda_2(\xi_1,\xi_2).
	\]
	由于 $\lambda_1\ne\lambda_2$，故 $(\xi_1,\xi_2)=0$.
\end{proof}



\begin{lemma}
	设 $(V,(\cdot,\cdot))$ 是 $n$ 维欧氏空间，$T\in\End(V)$. 则 $T\in\mathcal{S}(V)$ 当且仅当 $T$ 在 $V$ 的任意一组标准正交基下的矩阵都是对称阵.
\end{lemma}
\begin{proof}
	注意到 $T\in\mathcal{S}(V)$ 当且仅当对 $V$ 的任意基 $\alpha_1,\cdots,\alpha_n$ 都有
	\[
	(T(\alpha_i),\alpha_j)=(\alpha_i,T(\alpha_j)),
	\qquad 1\leq i,j\leq n.
	\]
	任取 $V$ 的一组标准正交基 $\varepsilon_1,\cdots,\varepsilon_n$，设
	\[
	(T(\varepsilon_1)\ \cdots\ T(\varepsilon_n))=(\varepsilon_1\ \cdots\ \varepsilon_n)A.
	\]
	则
	\[
	T(\varepsilon_j)=\sum_{k=1}^{n}A(k,j)\varepsilon_k.
	\]
	于是
	\[
	\begin{aligned}
		(T(\varepsilon_i),\varepsilon_j)
		&=\left(\sum_{k=1}^{n}A(k,i)\varepsilon_k,\varepsilon_j\right)
		=\sum_{k=1}^{n}A(k,i)(\varepsilon_k,\varepsilon_j)=A(j,i),\\
		(\varepsilon_i,T(\varepsilon_j))
		&=\left(\varepsilon_i,\sum_{k=1}^{n}A(k,j)\varepsilon_k\right)
		=\sum_{k=1}^{n}A(k,j)(\varepsilon_i,\varepsilon_k)=A(i,j).
	\end{aligned}
	\]
	故 $T\in\mathcal{S}(V)$ 当且仅当 $A(j,i)=A(i,j)$，即 $A'=A$. 反向证明即为以上每步可逆.
\end{proof}

\begin{example}
	设 $A\in M_n(\mathbb{R})$. 把 $A$ 看成 $\mathbb{R}^n$ 在标准内积下的线性变换
	\[
	X\longmapsto AX.
	\]
	则 $A\in\mathcal{S}(\mathbb{R}^n)$ 当且仅当 $A$ 是对称阵.
\end{example}
\begin{proof}
	法一：$\mathbb{R}^n$ 的基本向量 $e_1,\cdots,e_n$ 构成一组标准正交基，且 $A$ 在这组标准正交基下的矩阵就是 $A$ 本身. 由上一引理即得结论.
	
	法二：对任意 $X,Y\in\mathbb{R}^n$，有
	\[
	(A(X),Y)=(AX,Y)=(AX)'Y=X'A'Y,
	\]
	而
	\[
	(X,A(Y))=(X,AY)=X'(AY)=X'AY.
	\]
	所以 $(A(X),Y)=(X,A(Y))$ 对任意 $X,Y\in\mathbb{R}^n$ 都成立当且仅当 $A'=A$.
\end{proof}

\begin{corollary}
	设 $(V,(\cdot,\cdot))$ 是 $n$ 维欧氏空间. 则有线性空间同构
	\[
	\mathcal{S}(V)\cong \mathcal{S}_n(\mathbb{R})
	:=\{A\in M_n(\mathbb{R}):A'=A\}.
	\]
	特别地，
	\[
	\dim \mathcal{S}(V)=\frac12 n(n+1).
	\]
\end{corollary}
\begin{proof}
	任意取定 $V$ 的一组标准正交基 $\varepsilon_1,\cdots,\varepsilon_n$. 定义
	\[
	\rho:\mathcal{S}(V)\longrightarrow \mathcal{S}_n(\mathbb{R}),
	\qquad
	T\longmapsto T\text{ 在 }\varepsilon_1,\cdots,\varepsilon_n\text{ 下的矩阵 }A.
	\]
	由线性映射与矩阵的对应可知 $\rho$ 是线性映射且是单射. 只需证 $\rho$ 是满射.
	
	对任意 $A\in\mathcal{S}_n(\mathbb{R})$，定义 $T\in\End(V)$ 为
	\[
	(T(\varepsilon_1),\cdots,T(\varepsilon_n))=(\varepsilon_1\cdots\varepsilon_n)A.
	\]
	对任意
	\[
	\alpha=(\varepsilon_1\cdots\varepsilon_n)X,
	\qquad
	\beta=(\varepsilon_1\cdots\varepsilon_n)Y,
	\]
	有 $T(\alpha)$，$T(\beta)$ 的坐标分别为 $AX,AY$. 从而
	\[
	(T(\alpha),\beta)=(AX)'Y=X'A'Y=X'(AY)=(\alpha,T(\beta)).
	\]
	故 $T\in\mathcal{S}(V)$，即 $\rho$ 是满射. 因此 $\mathcal{S}(V)\cong\mathcal{S}_n(\mathbb{R})$. 又 $n$ 阶实对称阵由主对角线及其上方元素唯一确定，故维数为 $\dfrac12 n(n+1)$.
\end{proof}


\begin{lemma}
	设 $A\in M_n(\mathbb{R})$ 且 $A'=A$. 则 $A$ 的特征值全为实数.
\end{lemma}
\begin{proof}
	任取 $A$ 的一个特征值 $\lambda$. 设存在 $0\ne\xi\in\mathbb{C}^n$，使得
	\[
	A\xi=\lambda\xi.
	\]
	由于 $A\in M_n(\mathbb{R})$，取共轭得
	\[
	A\bar\xi=\bar\lambda\bar\xi.
	\]
	对上式取转置，并用 $A'=A$，得
	\[
	\bar\xi' A=\bar\lambda\bar\xi'.
	\]
	另一方面，由 $A\xi=\lambda\xi$ 有
	\[
	\bar\xi' A\xi=\lambda\bar\xi'\xi.
	\]
	而由 $\bar\xi' A=\bar\lambda\bar\xi'$ 有
	\[
	\bar\xi' A\xi=\bar\lambda\bar\xi'\xi.
	\]
	设 $\xi=(z_1,\cdots,z_n)'\in\mathbb{C}^n$，则
	\[
	\bar\xi'\xi=|z_1|^2+\cdots+|z_n|^2>0.
	\]
	故 $\lambda=\bar\lambda$，即 $\lambda\in\mathbb{R}$.
\end{proof}
$n$维欧式空间上的任意对称变换都有$n$个特征值
\begin{example}
	设 $C\in M_n(\mathbb{R})$ 且 $C'=-C$. 则 $C$ 的特征值只能是 $0$ 或纯虚数.
\end{example}
\begin{proof}
	设 $C\xi=\lambda\xi$，$0\ne\xi\in\mathbb{C}^n$. 与上面类似可得
	\[
	\bar\xi'C=-\bar\lambda\bar\xi'.
	\]
	于是
	\[
	\lambda\bar\xi'\xi=\bar\xi'C\xi=-\bar\lambda\bar\xi'\xi.
	\]
	由于 $\bar\xi'\xi>0$，故 $\lambda=-\bar\lambda$. 因而 $\lambda$ 的实部为 $0$，即 $\lambda=0$ 或 $\lambda$ 为纯虚数.
\end{proof}

\begin{theorem}
	设 $A\in M_n(\mathbb{R})$ 且 $A'=A$. 则 $A$ 在实数域上可对角化. 进一步，$A$ 正交相似于对角阵.
\end{theorem}
\begin{proof}
	法一：验证 $\mathbb{R}^n$ 可以分解为 $A$ 的特征子空间的直和.
	
	设 $\lambda_1,\cdots,\lambda_s$ 是 $A$ 的全部互不相同的特征值. 由上面的引理可知 $\lambda_i\in\mathbb{R}$. 设
	\[
	V_{\lambda_i}=\{X\in\mathbb{R}^n:AX=\lambda_iX\},
	\qquad 1\leq i\leq s.
	\]
	令
	\[
	W=V_{\lambda_1}\oplus\cdots\oplus V_{\lambda_s}.
	\]
	下证 $W=\mathbb{R}^n$. 若不然，则 $W$ 是 $\mathbb{R}^n$ 的真子空间. 考虑 $\mathbb{R}^n$ 上的标准内积，则 $W^\perp$ 是欧氏空间. 把 $A$ 看成欧氏空间 $(\mathbb{R}^n,(\cdot,\cdot))$ 上的线性变换，由于 $A$ 是对称变换且 $W$ 是 $A$-子空间，所以 $W^\perp$ 也是 $A$-子空间. 从而 $A|_{W^\perp}$ 是 $W^\perp$ 上的对称变换.
	
	由上面的引理，$A|_{W^\perp}$ 的特征值也全为实数. 任取 $A|_{W^\perp}$ 的一个特征值 $\mu$，则 $\mu$ 也是 $A$ 的特征值. 又对应的非零特征向量属于 $W^\perp$，不属于 $W$，故 $\mu\ne\lambda_i$，$1\leq i\leq s$. 这与 $\lambda_1,\cdots,\lambda_s$ 是 $A$ 的全部互不相同的特征值矛盾. 因此
	\[
	\mathbb{R}^n=W=V_{\lambda_1}\oplus\cdots\oplus V_{\lambda_s}.
	\]
	取 $V_{\lambda_i}$ 的一组标准正交基
	\[
	\xi_{i1},\cdots,\xi_{ir_i},
	\qquad 1\leq i\leq s,
	\]
	则
	\[
	\xi_{11},\cdots,\xi_{1r_1},\xi_{21},\cdots,\xi_{2r_2},\cdots,\xi_{s1},\cdots,\xi_{sr_s}
	\]
	是 $\mathbb{R}^n$ 的一组标准正交基. 令
	\[
	U=(\xi_{11}\cdots\xi_{1r_1}\ \xi_{21}\cdots\xi_{2r_2}\ \cdots\ \xi_{s1}\cdots\xi_{sr_s}).
	\]
	则 $U$ 是正交阵，且
	\[
	U^{-1}AU=\operatorname{diag}(\lambda_1E_{r_1},\lambda_2E_{r_2},\cdots,\lambda_sE_{r_s}).
	\]
	
	法二：对 $n$ 作归纳，证明存在正交阵 $U$ 使得 $U^{-1}AU$ 是对角阵.
	
	当 $n=1$ 时，结论显然成立，取 $U=E_1=(1)$ 即可. 设结论对任意 $n-1$ 阶实对称阵都成立. 对于 $A\in M_n(\mathbb{R})$，$A'=A$. 由上面的引理，可以任取 $A$ 的一个特征值 $\lambda\in\mathbb{R}$. 取 $\xi\in\mathbb{R}^n$ 满足 $|\xi|=1$ 且
	\[
	A\xi=\lambda\xi.
	\]
	把 $\xi$ 扩充为 $\mathbb{R}^n$ 的一组标准正交基
	\[
	\xi,\xi_2,\cdots,\xi_n.
	\]
	设
	\[
	U_1=(\xi\ \xi_2\ \cdots\ \xi_n).
	\]
	则 $U_1\in O_n(\mathbb{R})$，且
	\[
	U_1^{-1}AU_1=\begin{pmatrix}
		\lambda&\delta'\\
		0&A_1
	\end{pmatrix},
	\qquad \delta\in\mathbb{R}^{n-1}.
	\]
	两边取转置，并用 $U_1^{-1}=U_1'$ 与 $A'=A$，得
	\[
	\begin{pmatrix}
		\lambda&0\\
		\delta&A_1'
	\end{pmatrix}
	=\begin{pmatrix}
		\lambda&\delta'\\
		0&A_1
	\end{pmatrix}.
	\]
	故 $\delta=0$ 且 $A_1'=A_1$. 因此
	\[
	U_1^{-1}AU_1=\begin{pmatrix}
		\lambda&0\\
		0&A_1
	\end{pmatrix},
	\qquad A_1\in M_{n-1}(\mathbb{R}),\quad A_1'=A_1.
	\]
	由归纳假设，存在 $U_2\in O_{n-1}(\mathbb{R})$，使得 $U_2'A_1U_2=D$ 是对角阵. 令
	\[
	U=U_1\begin{pmatrix}
		1&0\\
		0&U_2
	\end{pmatrix}\in O_n(\mathbb{R}).
	\]
	于是
	\[
	\begin{aligned}
		U^{-1}AU
		&=\begin{pmatrix}1&0\\0&U_2\end{pmatrix}^{-1}
		U_1^{-1}AU_1
		\begin{pmatrix}1&0\\0&U_2\end{pmatrix}\\
		&=\begin{pmatrix}1&0\\0&U_2^{-1}\end{pmatrix}
		\begin{pmatrix}\lambda&0\\0&A_1\end{pmatrix}
		\begin{pmatrix}1&0\\0&U_2\end{pmatrix}
		=\begin{pmatrix}\lambda&0\\0&D\end{pmatrix}.
	\end{aligned}
	\]
	故 $A$ 正交相似于对角阵.
\end{proof}

\begin{corollary}
	设 $T$ 是欧氏空间 $V$ 的对称变换. 则存在 $V$ 的一组标准正交基，使得 $T$ 在这组基下的矩阵是对角阵.
\end{corollary}

\begin{example}
	设 $A$ 是实对称阵，且 $A$ 的特征多项式为
	\[
	f(\lambda)=(\lambda-2)^3(\lambda+3)^2.
	\]
	求 $A$ 的极小多项式与初等因子.
\end{example}
\begin{proof}[解答]
	由于 $A$ 在 $\mathbb{R}$ 上可对角化，且
	\[
	A\sim\begin{pmatrix}
		2E_3&0\\
		0&-3E_2
	\end{pmatrix},
	\]
	所以 $A$ 的极小多项式为
	\[
	(\lambda-2)(\lambda+3).
	\]
	$A$ 在 $\mathbb{R}$ 上的初等因子为
	\[
	\lambda-2,\quad \lambda-2,\quad \lambda-2,
	\qquad
	\lambda+3,
	\quad
	\lambda+3.
	\]
\end{proof}

\begin{example}
	设 $A,B\in M_n(\mathbb{R})$，且 $A'=A$，$B'=B$. 则
	\[
	A\sim B\Longleftrightarrow A,B\text{ 的特征多项式相同}.
	\]
\end{example}
\begin{proof}
	必要性显然. 反之，若 $A,B$ 的特征多项式相同，则二者的特征值及其代数重数完全相同. 又 $A,B$ 均为实对称阵，故都相似于同一个由特征值按重数排列而成的对角阵. 因此 $A\sim B$.
\end{proof}

\begin{example}
	设 $n$ 阶方阵 $A$ 的元素全为 $1$. 判断 $A$ 在 $\mathbb{R}$ 上是否可对角化；如果可对角化，求与 $A$ 相似的对角阵.
\end{example}
\begin{proof}[解答]
	因为 $A$ 是实对称阵，所以 $A$ 在 $\mathbb{R}$ 上可对角化. 由 $|A|=0$ 可知 $0$ 是 $A$ 的一个特征值，且 $A$ 的秩为 $1$，所以 $0$ 的几何重数为 $n-1$. 因为 $A$ 可对角化，所以 $0$ 的代数重数也为 $n-1$.
	
	于是 $A$ 有且仅有一个非零特征值 $\lambda$，且 $\lambda$ 的重数为 $1$. 由
	\[
	n=\operatorname{tr}(A)=(n-1)\cdot0+\lambda
	\]
	可得 $\lambda=n$. 故
	\[
	A\sim\begin{pmatrix}
		0_{n-1}&0\\
		0&n
	\end{pmatrix}.
	\]
\end{proof}

\begin{example}
	设 $A\in M_n(\mathbb{R})$，$A'=A$. 则
	\[
	A=0\Longleftrightarrow A^2=0.
	\]
\end{example}
\begin{proof}
	只需证 $A^2=0\Rightarrow A=0$. 由 $A'=A$，存在正交阵 $U$ 使得
	\[
	U^{-1}AU=\operatorname{diag}(\lambda_1,\cdots,\lambda_n).
	\]
	于是
	\[
	U^{-1}A^2U=\operatorname{diag}(\lambda_1^2,\cdots,\lambda_n^2)=0.
	\]
	故 $\lambda_i=0$，$1\leq i\leq n$，从而 $A=0$.
\end{proof}

\begin{example}
	设 $A\in M_n(\mathbb{R})$，$A'=A$. 则 $A$ 的特征值全大于 $0$ 当且仅当存在可逆的 $B\in M_n(\mathbb{R})$，$B'=B$，使得
	\[
	A=B^2.
	\]
\end{example}
\begin{proof}
	若 $A$ 的特征值全大于 $0$，取正交阵 $U$ 使
	\[
	U^{-1}AU=\operatorname{diag}(\lambda_1,\cdots,\lambda_n),
	\qquad \lambda_i>0.
	\]
	令
	\[
	B=U\operatorname{diag}(\sqrt{\lambda_1},\cdots,\sqrt{\lambda_n})U^{-1}.
	\]
	则 $B'=B$，$B$ 可逆，且 $B^2=A$.
	
	反之，若 $A=B^2$，其中 $B'=B$ 且 $B$ 可逆，则 $B$ 正交相似于实对角阵，故 $A=B^2$ 的特征值为 $B$ 的非零特征值的平方，均大于 $0$.
\end{proof}

\begin{example}
	设 $A,B\in M_n(\mathbb{R})$，$A'=A$，$B'=B$. 若
	\[
	X'AX=X'BX
	\]
	对任意 $X\in\mathbb{R}^n$ 都成立，则 $A=B$.
\end{example}
\begin{proof}
	令 $C=A-B$，则 $C'=C$ 且 $X'CX=0$ 对任意 $X\in\mathbb{R}^n$ 成立. 取 $X=e_i$ 得 $c_{ii}=0$；取 $X=e_i+e_j$，$i\ne j$，得
	\[
	0=(e_i+e_j)'C(e_i+e_j)=c_{ii}+c_{ij}+c_{ji}+c_{jj}=2c_{ij}.
	\]
	故 $c_{ij}=0$，从而 $C=0$，即 $A=B$.
\end{proof}


\begin{algo}(实对称阵正交对角化的算法)
	给定 $n$ 阶实对称阵 $A$，求正交阵 $U$ 使得 $U^{-1}AU$ 是对角阵，可以按如下步骤进行：
	\begin{enumerate}
		\item 求出 $A$ 的全部互不相同的特征值 $\lambda_1,\cdots,\lambda_s$；
		\item 对每个 $\lambda_i$，求方程组
		\[
		(\lambda_iE_n-A)X=0
		\]
		的一个基础解系 $\eta_{i1},\cdots,\eta_{ir_i}$；
		\item 对 $\eta_{i1},\cdots,\eta_{ir_i}$ 作 Schmidt 正交化，再单位化，得到
		\[
		\xi_{i1},\cdots,\xi_{ir_i};
		\]
		\item 令
		\[
		U=(\xi_{11}\cdots\xi_{1r_1}\ \cdots\ \xi_{s1}\cdots\xi_{sr_s}).
		\]
		则 $U$ 是正交阵，且
		\[
		U^{-1}AU=U'AU=\operatorname{diag}(\lambda_1E_{r_1},\cdots,\lambda_sE_{r_s}).
		\]
		\end{enumerate}
\end{algo}

\begin{example}
	设
	\[
	A=\begin{pmatrix}
		2&2&-2\\
		2&5&-4\\
		-2&-4&5
	\end{pmatrix}.
	\]
	求正交阵 $P$ 使得 $P^{-1}AP$ 是对角阵.
\end{example}
\begin{proof}[解答]
	先求特征多项式：
	\[
	\begin{aligned}
		f(\lambda)
		&=\left|\lambda E_3-A\right|\\
		&=\begin{vmatrix}
			\lambda-2&-2&2\\
			-2&\lambda-5&4\\
			2&4&\lambda-5
		\end{vmatrix}
		=(\lambda-1)^2(\lambda-10).
	\end{aligned}
	\]
	因此 $A$ 的全部互不相同的特征值为 $\lambda_1=1$，$\lambda_2=10$.
	
	方程组 $(\lambda_1E_3-A)X=0$ 的一个基础解系为
	\[
	\eta_{11}=\begin{pmatrix}2\\0\\1\end{pmatrix},
	\qquad
	\eta_{12}=\begin{pmatrix}-2\\1\\0\end{pmatrix}.
	\]
	作 Schmidt 正交化，再单位化得
	\[
	\xi_{11}=\frac1{\sqrt5}\begin{pmatrix}2\\0\\1\end{pmatrix},
	\qquad
	\xi_{12}=\frac{\sqrt5}{3}\begin{pmatrix}-\dfrac25\\1\\\dfrac45\end{pmatrix}.
	\]
	方程组 $(\lambda_2E_3-A)X=0$ 的一个基础解系为
	\[
	\eta_{21}=\begin{pmatrix}-1\\-2\\2\end{pmatrix},
	\]
	单位化得
	\[
	\xi_{21}=\frac13\begin{pmatrix}-1\\-2\\2\end{pmatrix}.
	\]
	令
	\[
	P=(\xi_{11}\ \xi_{12}\ \xi_{21}),
	\]
	则 $P$ 是正交阵，且
	\[
	P^{-1}AP=\begin{pmatrix}
		1&0&0\\
		0&1&0\\
		0&0&10
	\end{pmatrix}.
	\]
\end{proof}

\section{伴随算子、欧氏空间上的正规算子}



\begin{lemma}
	设 $(V,(\cdot,\cdot))$ 是实内积空间，$T\in\End(V)$. 如果 $S_1,S_2\in\End(V)$ 满足：对任意 $\alpha,\beta\in V$ 都有
	\[
	(T(\alpha),\beta)=(\alpha,S_1(\beta))=(\alpha,S_2(\beta)),
	\]
	则 $S_1=S_2$.
\end{lemma}
\begin{proof}
	对任意 $\beta,\alpha\in V$，由
	\[
	(\alpha,S_1(\beta))=(\alpha,S_2(\beta))
	\]
	可知
	\[
	(\alpha,S_1(\beta)-S_2(\beta))=0.
	\]
	由内积的正定性推出 $S_1(\beta)=S_2(\beta)$ 对任意 $\beta\in V$ 都成立，因此 $S_1=S_2$.
\end{proof}

\begin{definition}
	设 $(V,(\cdot,\cdot))$ 是实内积空间，$T\in\End(V)$. 如果 $S\in\End(V)$ 满足：对任意 $\alpha,\beta\in V$ 都有
	\[
	(T(\alpha),\beta)=(\alpha,S(\beta)),
	\]
	则称 $S$ 是 $T$ 的\textbf{伴随变换}（伴随算子，adjoint operator）.
\end{definition}

由上面的引理，如果 $T$ 的伴随算子存在，则必然唯一. 下面说明在欧氏空间中，任意线性变换都有伴随算子.

\begin{lemma}
	设 $(V,(\cdot,\cdot))$ 是实内积空间. 对任意 $\gamma\in V$，定义
	\[
	f_\gamma:V\longrightarrow\mathbb{R},
	\qquad
	f_\gamma(\gamma')=(\gamma,\gamma'),
	\qquad \gamma'\in V.
	\]
	则 $f_\gamma\in V^*$. 进一步，映射
	\[
	\rho:V\longrightarrow V^*:=\Hom(V,\mathbb{R}),
	\qquad
	\gamma\longmapsto f_\gamma
	\]
	是单的线性映射.
\end{lemma}
\begin{proof}
	首先，$f_\gamma$ 关于第二个变量线性，故 $f_\gamma\in V^*$. 对任意 $k_1,k_2\in\mathbb{R}$ 与 $\gamma_1,\gamma_2\in V$，有
	\[
	f_{k_1\gamma_1+k_2\gamma_2}=k_1f_{\gamma_1}+k_2f_{\gamma_2},
	\]
	故 $\rho$ 是线性映射.
	
	又
	\[
	\gamma\in\ker\rho
	\Longleftrightarrow f_\gamma=0
	\Longleftrightarrow (\gamma,\gamma')=0\text{ 对任意 }\gamma'\in V\text{ 都成立}
	\Longleftrightarrow \gamma=0.
	\]
	因此 $\ker\rho=0$，即 $\rho$ 是单射.
\end{proof}

\begin{corollary}
	设 $(V,(\cdot,\cdot))$ 是欧氏空间. 则上述
	\[
	\rho:V\longrightarrow V^*
	\]
	是同构映射.
\end{corollary}
\begin{proof}
	设 $\dim V=n$，则 $\dim V^*=n$. 由
	\[
	\dim\ker\rho+\dim\Im \rho=\dim V=n=\dim V^*
	\]
	和 $\ker\rho=0$ 可知 $\rho$ 是满射. 因此 $\rho$ 是同构映射.
\end{proof}

\begin{corollary}
	设 $(V,(\cdot,\cdot))$ 是欧氏空间. 则任意 $T\in\End(V)$ 都有伴随算子，且伴随算子必然唯一.
\end{corollary}
\begin{proof}
	任意固定 $\beta\in V$. 于是有 $f_\beta\in V^*$，从而由 $T\in\End(V)$ 可得复合映射
	\[
	f_\beta T\in V^*.
	\]
	由上面的同构映射 $\rho:V\to V^*$ 可知，存在唯一的 $\beta'\in V$，使得
	\[
	\rho(\beta')=f_\beta T.
	\]
	从而有映射
	\[
	T^*:V\longrightarrow V,
	\qquad
	\beta\longmapsto \beta'
	\]
	使得 $\rho(T^*(\beta))=f_\beta T$. 即对任意 $\alpha\in V$，有
	\[
	\begin{aligned}
		(f_\beta T)(\alpha)
		&=f_\beta(T(\alpha))=(\beta,T(\alpha))\\
		&=\rho(T^*(\beta))(\alpha)=f_{T^*(\beta)}(\alpha)\\
		&=(T^*(\beta),\alpha)=(\alpha,T^*(\beta)).
	\end{aligned}
	\]
	所以 $T^*$ 是 $T$ 的伴随算子.
	
	下面验证 $T^*$ 是线性的. 对任意 $\beta_1,\beta_2\in V$，有
	\[
	\begin{aligned}
		\rho(T^*(\beta_1+\beta_2))
		&=f_{\beta_1+\beta_2}T\\
		&=f_{\beta_1}T+f_{\beta_2}T\\
		&=\rho(T^*(\beta_1))+\rho(T^*(\beta_2))\\
		&=\rho(T^*(\beta_1)+T^*(\beta_2)).
	\end{aligned}
	\]
	由于 $\rho$ 是同构，故
	\[
	T^*(\beta_1+\beta_2)=T^*(\beta_1)+T^*(\beta_2).
	\]
	类似地可验证 $T^*$ 保数乘. 因而 $T^*\in\End(V)$.
	
	伴随算子的唯一性由前面的引理给出.
\end{proof}

记号：用 $T^*$ 表示 $T$ 的伴随算子.

\begin{example}
	设 $(V,(\cdot,\cdot))$ 是欧氏空间.
	\begin{enumerate}
		\item 如果 $T$ 是 $V$ 上的正交变换，则
		\[
		T^*=T^{-1}.
		\]
		\item 如果 $T$ 是 $V$ 上的对称变换，则
		\[
		T^*=T.
		\]
	\end{enumerate}
\end{example}
\begin{proof}
	若 $T$ 是正交变换，则对任意 $\alpha,\beta\in V$，有
	\[
	(T(\alpha),\beta)=(T(\alpha),T(T^{-1}(\beta)))=(\alpha,T^{-1}(\beta)),
	\]
	故 $T^*=T^{-1}$.
	
	若 $T$ 是对称变换，则对任意 $\alpha,\beta\in V$，有
	\[
	(T(\alpha),\beta)=(\alpha,T(\beta)),
	\]
	故 $T^*=T$.
\end{proof}

\begin{corollary}
	设 $(V,(\cdot,\cdot))$ 是 $n$ 维欧氏空间，$T\in\End(V)$. 设 $T$ 在 $V$ 的一个标准正交基
	\[
	\varepsilon_1,\cdots,\varepsilon_n
	\]
	下的矩阵为 $A$. 则 $T$ 的伴随算子 $T^*$ 在 $\varepsilon_1,\cdots,\varepsilon_n$ 下的矩阵为 $A'$.
\end{corollary}
\begin{proof}
	设 $T^*$ 在 $\varepsilon_1,\cdots,\varepsilon_n$ 下的矩阵为 $B$. 则
	\[
	T^*(\varepsilon_i)=\sum_{k=1}^{n}B(k,i)\varepsilon_k,
	\qquad
	T(\varepsilon_j)=\sum_{\ell=1}^{n}A(\ell,j)\varepsilon_\ell.
	\]
	于是
	\[
	\begin{aligned}
		(T^*(\varepsilon_i),\varepsilon_j)
		&=\left(\sum_{k=1}^{n}B(k,i)\varepsilon_k,\varepsilon_j\right)=B(j,i),\\
		(\varepsilon_i,T(\varepsilon_j))
		&=\left(\varepsilon_i,\sum_{\ell=1}^{n}A(\ell,j)\varepsilon_\ell\right)=A(i,j).
	\end{aligned}
	\]
	又由伴随算子的定义，
	\[
	(T^*(\varepsilon_i),\varepsilon_j)=(\varepsilon_i,T(\varepsilon_j)),
	\]
	故 $B(j,i)=A(i,j)$，从而 $B=A'$.
\end{proof}

\begin{corollary}
	设 $(V,(\cdot,\cdot))$ 是欧氏空间，$T_1,T_2\in\End(V)$. 则
	\[
	(T_1+T_2)^*=T_1^*+T_2^*,
	\qquad
	(T_1T_2)^*=T_2^*T_1^*.
	\]
\end{corollary}
\begin{proof}
	取定 $V$ 的一个标准正交基，利用线性变换与矩阵之间的一一对应关系. 若 $T_1,T_2$ 在该基下的矩阵分别为 $A_1,A_2$，则由上一推论，$T_1^*,T_2^*$ 的矩阵分别为 $A_1',A_2'$. 因为
	\[
	(A_1+A_2)'=A_1'+A_2',
	\qquad
	(A_1A_2)'=A_2'A_1',
	\]
	故
	\[
	(T_1+T_2)^*=T_1^*+T_2^*,
	\qquad
	(T_1T_2)^*=T_2^*T_1^*.
	\]
\end{proof}

\begin{definition}
	设 $(V,(\cdot,\cdot))$ 是实内积空间，$T\in\End(V)$. 如果 $T$ 的伴随算子 $T^*$ 存在且
	\[
	TT^*=T^*T,
	\]
	则称 $T$ 是\textbf{正规变换}（正规算子，normal operator）.
\end{definition}

\begin{example}
	欧氏空间上的正交变换和对称变换都是正规算子.
\end{example}
\begin{proof}
	若 $T$ 是正交变换，则 $T^*=T^{-1}$，于是
	\[
	TT^*=TT^{-1}=\id_V=T^{-1}T=T^*T.
	\]
	若 $T$ 是对称变换，则 $T^*=T$，于是
	\[
	TT^*=T^2=T^*T.
	\]
	故二者都是正规算子.
\end{proof}

\begin{definition}
	设 $A\in M_n(\mathbb{R})$. 如果
	\[
	AA'=A'A,
	\]
	则称 $A$ 是一个\textbf{正规矩阵}.
\end{definition}

\begin{example}
	正交阵和实对称阵都是正规矩阵.
\end{example}
\begin{proof}
	若 $A$ 是正交阵，则 $A'A=E_n=AA'$. 若 $A$ 是实对称阵，则 $A'=A$，于是
	\[
	AA'=A^2=A'A.
	\]
	故正交阵和实对称阵都是正规矩阵.
\end{proof}




\begin{lemma}
	设 $A\in M_2(\mathbb{R})$. 若 $A$ 是正规矩阵，则 $A$ 正交相似于下列形式之一：
	\[
	\begin{pmatrix}
		a_1&0\\
		0&a_2
	\end{pmatrix},
	\qquad
	\begin{pmatrix}
		a&b\\
		-b&a
	\end{pmatrix},
	\]
	其中 $a_1,a_2,a,b\in\mathbb{R}$.
\end{lemma}

\begin{proof}
	设
	\[
	A=
	\begin{pmatrix}
		x&y\\
		z&t
	\end{pmatrix}.
	\]
	由于 $A$ 是正规矩阵，即
	\[
	AA'=A'A.
	\]
	直接计算得
	\[
	AA'=
	\begin{pmatrix}
		x^2+y^2&xz+yt\\
		xz+yt&z^2+t^2
	\end{pmatrix},
	\qquad
	A'A=
	\begin{pmatrix}
		x^2+z^2&xy+zt\\
		xy+zt&y^2+t^2
	\end{pmatrix}.
	\]
	比较 $(1,1)$ 元可得
	\[
	y^2=z^2.
	\]
	故
	\[
	z=y \quad \text{或} \quad z=-y.
	\]
	
	若 $z=y$，则
	\[
	A=
	\begin{pmatrix}
		x&y\\
		y&t
	\end{pmatrix}
	\]
	是实对称矩阵. 因此 $A$ 正交相似于对角阵，即存在正交矩阵 $U$，使得
	\[
	U^{-1}AU=
	\begin{pmatrix}
		a_1&0\\
		0&a_2
	\end{pmatrix}.
	\]
	
	若 $z=-y$，则
	\[
	A=
	\begin{pmatrix}
		x&y\\
		-y&t
	\end{pmatrix}.
	\]
	再比较 $AA'=A'A$ 的 $(1,2)$ 元，得
	\[
	xz+yt=xy+zt.
	\]
	代入 $z=-y$，有
	\[
	-xy+yt=xy-yt,
	\]
	即
	\[
	2y(t-x)=0.
	\]
	若 $y=0$，则 $A$ 已是对角阵. 若 $y\ne0$，则 $x=t$. 于是
	\[
	A=
	\begin{pmatrix}
		x&y\\
		-y&x
	\end{pmatrix}.
	\]
	令 $a=x,\ b=y$，即得
	\[
	A=
	\begin{pmatrix}
		a&b\\
		-b&a
	\end{pmatrix}.
	\]
	
	综上，$A$ 正交相似于
	\[
	\begin{pmatrix}
		a_1&0\\
		0&a_2
	\end{pmatrix}
	\quad \text{或} \quad
	\begin{pmatrix}
		a&b\\
		-b&a
	\end{pmatrix}.
	\]
\end{proof}
\begin{theorem}
	设 $A\in M_n(\mathbb{R})$. 设 $A$ 是正规矩阵. 则 $A$ 正交相似于如下形式的准对角阵：
	\[
	\operatorname{diag}\left(
	(a_1),(a_2),\cdots,(a_m),
	\begin{pmatrix}b_1&c_1\\-c_1&b_1\end{pmatrix},
	\cdots,
	\begin{pmatrix}b_s&c_s\\-c_s&b_s\end{pmatrix}
	\right),
	\]
	其中 $a_i,b_j,c_j\in\mathbb{R}$.
\end{theorem}


\section{酉空间简介}

欧氏空间是建立在实数域 $\mathbb{R}$ 上的线性空间上的，它在很多地方都很有用.
现在我们希望把它推广到复数域 $\mathbb{C}$ 上的线性空间.
但是，原来实内积空间中的内积定义在复数域上不再适用. 因为若
\[
(\alpha,\alpha)>0,
\]
则
\[
(i\alpha,i\alpha)=i^2(\alpha,\alpha)<0.
\]
因此，有必要对复线性空间上的内积定义作出修正.

\begin{definition}
	设 $V$ 是复数域 $\mathbb{C}$ 上的线性空间. $V$ 上的一个内积是指一个映射
	\[
	V\times V\longrightarrow \mathbb{C},
	\qquad
	\{\alpha,\beta\}\longmapsto (\alpha,\beta),
	\]
	满足如下条件：
	\begin{enumerate}
		\item $(\alpha,\beta)=\overline{(\beta,\alpha)}$;
		\item $(k_1\alpha_1+k_2\alpha_2,\beta)
		=k_1(\alpha_1,\beta)+k_2(\alpha_2,\beta)$,
		其中 $k_1,k_2\in\mathbb{C}$, $\alpha_1,\alpha_2,\beta\in V$;
		\item $(\alpha,\alpha)>0$, 若 $\alpha\ne0$.
	\end{enumerate}
	配备了内积的有限维复线性空间称为酉空间.
\end{definition}

由定义可以直接推出
\[
(\alpha,0)=0,
\qquad
\forall \alpha\in V,
\]
并且
\[
(\alpha,k_1\beta_1+k_2\beta_2)
=\overline{k_1}(\alpha,\beta_1)+\overline{k_2}(\alpha,\beta_2).
\]

\begin{example}
	令 $V=\mathbb{C}^n$. 对
	\[
	\alpha=(a_1,\cdots,a_n),
	\qquad
	\beta=(b_1,\cdots,b_n),
	\]
	定义
	\[
	(\alpha,\beta)=\alpha\overline{\beta}^{\,T}
	=a_1\overline{b_1}+\cdots+a_n\overline{b_n}.
	\]
	则 $\mathbb{C}^n$ 是一个酉空间.
\end{example}

酉空间有许多与欧氏空间类似的性质.

\begin{enumerate}
	\item 可以定义向量 $\alpha$ 的长度为
	\[
	|\alpha|=\sqrt{(\alpha,\alpha)}.
	\]
	
	\item Cauchy--Bunyakovski 不等式成立：
	\[
	|(\alpha,\beta)|\leq |\alpha|\,|\beta|.
	\]
	等号成立当且仅当 $\alpha,\beta$ 线性相关.
	
	\item 若
	\[
	(\alpha,\beta)=0,
	\]
	则称 $\alpha$ 与 $\beta$ 正交，记作
	\[
	\alpha\perp\beta.
	\]
	
	\item 可以用显然的方式定义标准正交基，并且同样有 Schmidt 正交化过程.
	
	\item 若 $W$ 是 $V$ 的子空间，则可以定义 $W$ 的正交补 $W^\perp$, 并且有分解
	\[
	V=W\oplus W^\perp.
	\]
	
	\item 设 $A$ 是复矩阵. 记 $\overline{A}$ 为 $A$ 的每个元素取共轭后所得的矩阵.
	若
	\[
	\overline{A}^{\,T}A=A\overline{A}^{\,T}=E,
	\]
	则称 $A$ 是酉矩阵. 对酉矩阵 $A$, 有
	\[
	|\det A|=1.
	\]
	因此，酉矩阵一定可逆.
	
	从一个标准正交基到另一个标准正交基的过渡矩阵是酉矩阵.
	
	\item 设 $\mathcal{A}$ 是 $V$ 上的线性变换. 若
	\[
	(\mathcal{A}\alpha,\mathcal{A}\beta)=(\alpha,\beta),
	\qquad \alpha,\beta\in V,
	\]
	则称 $\mathcal{A}$ 是酉变换.
	
	酉变换 $\mathcal{A}$ 在标准正交基下的矩阵是酉矩阵.
	
	\item 设 $\mathcal{A}$ 是 $V$ 上的线性变换. 若
	\[
	(\mathcal{A}\alpha,\beta)=(\alpha,\mathcal{A}\beta),
	\qquad \forall \alpha,\beta\in V,
	\]
	则称 $\mathcal{A}$ 是对称变换.
	
	若复矩阵 $A$ 满足
	\[
	\overline{A}^{\,T}=A,
	\]
	则称 $A$ 是 Hermite 矩阵.
	
	对称变换在标准正交基下的矩阵是 Hermite 矩阵.
	
	\item Hermite 矩阵 $A$ 的特征值全为实数. 若 $\lambda_1\ne \lambda_2$ 是 $A$ 的特征值,
	且 $\xi_i$ 是属于 $\lambda_i$ 的特征向量 $(i=1,2)$, 则
	\[
	\xi_1\perp\xi_2.
	\]
	
	\item 对任意 Hermite 矩阵 $A$, 存在酉矩阵 $C$, 使得
	\[
	C^{-1}AC=\overline{C}^{\,T}AC=\operatorname{diag}\{d_1,\cdots,d_n\}.
	\]
	其中 $d_i,\ 1\leq i\leq n$, 是 $A$ 的特征值.
\end{enumerate}
\section{有价值的习题}

\begin{problem}
	设 $V$ 是欧氏空间，定义 $\alpha$ 与 $\beta$ 的距离为
	\[
	d(\alpha,\beta)=|\alpha-\beta|.
	\]
	证明：
	\begin{enumerate}
		\item $d(\alpha,\beta)\leqslant d(\alpha,\gamma)+d(\gamma,\beta)$，$\alpha,\beta,\gamma\in V$；
		\item 若 $\alpha,\beta\ne0$，则
		\[
		d(\alpha,\beta)
		=\sqrt{|\alpha|^2+|\beta|^2-2|\alpha|\,|\beta|\cos\angle(\alpha,\beta)}.
		\]
	\end{enumerate}
\end{problem}
\begin{solution}
	\begin{enumerate}
		\item 由 Cauchy--Bunyakovsky--Schwarz 不等式，
		\[
		\begin{aligned}
			(|\alpha-\gamma|+|\gamma-\beta|)^2
			&=|\alpha-\gamma|^2+|\gamma-\beta|^2
			+2|\alpha-\gamma|\,|\gamma-\beta| \\
			&\geqslant |\alpha-\gamma|^2+|\gamma-\beta|^2
			+2(\alpha-\gamma,\gamma-\beta) \\
			&=((\alpha-\gamma)+(\gamma-\beta),(\alpha-\gamma)+(\gamma-\beta)) \\
			&=|\alpha-\beta|^2.
		\end{aligned}
		\]
		两边开方即得结论.
		\item 由夹角定义
		\[
		\cos\angle(\alpha,\beta)=\frac{(\alpha,\beta)}{|\alpha|\,|\beta|},
		\]
		于是
		\[
		\begin{aligned}
			d(\alpha,\beta)^2
			&=|\alpha-\beta|^2
			=(\alpha-\beta,\alpha-\beta) \\
			&=|\alpha|^2+|\beta|^2-2(\alpha,\beta) \\
			&=|\alpha|^2+|\beta|^2-2|\alpha|\,|\beta|\cos\angle(\alpha,\beta).
		\end{aligned}
		\]
		开方即得所证.
	\end{enumerate}
\end{solution}

\begin{problem}

	证明任意欧氏空间 $V$ 中的平行四边形公式：对任意 $\alpha,\beta\in V$，
	\[
	|\alpha+\beta|^2+|\alpha-\beta|^2=2|\alpha|^2+2|\beta|^2.
	\]
\end{problem}
\begin{solution}
	直接展开得
	\[
	\begin{aligned}
		|\alpha+\beta|^2+|\alpha-\beta|^2
		&=(\alpha+\beta,\alpha+\beta)+(\alpha-\beta,\alpha-\beta) \\
		&=(\alpha,\alpha)+2(\alpha,\beta)+(\beta,\beta)
		+(\alpha,\alpha)-2(\alpha,\beta)+(\beta,\beta) \\
		&=2|\alpha|^2+2|\beta|^2.
	\end{aligned}
	\]
\end{solution}

\begin{problem}

	设 $\varepsilon_1,\cdots,\varepsilon_n$ 是欧氏空间 $V$ 的一组标准正交基. 证明：
	\begin{enumerate}
		\item 对任意 $\alpha\in V$，有
		\[
		\alpha=(\alpha,\varepsilon_1)\varepsilon_1+\cdots+(\alpha,\varepsilon_n)\varepsilon_n;
		\]
		\item 对任意 $\alpha\in V$，有
		\[
		|\alpha|^2=(\alpha,\varepsilon_1)^2+\cdots+(\alpha,\varepsilon_n)^2.
		\]
	\end{enumerate}
\end{problem}
\begin{solution}
	\begin{enumerate}
		\item 设
		\[
		\alpha=\sum_{i=1}^n x_i\varepsilon_i,
		\qquad x_i\in\R.
		\]
		对每个 $1\leqslant i\leqslant n$，两边与 $\varepsilon_i$ 作内积，得
		\[
		(\alpha,\varepsilon_i)=x_i(\varepsilon_i,\varepsilon_i)=x_i.
		\]
		代回坐标表达式即得结论.
		\item 由 (1) 与 $(\varepsilon_i,\varepsilon_j)=\delta_{ij}$，有
		\[
		\begin{aligned}
			|\alpha|^2
			&=(\alpha,\alpha)
			=\left(\sum_{i=1}^n(\alpha,\varepsilon_i)\varepsilon_i,
			\sum_{j=1}^n(\alpha,\varepsilon_j)\varepsilon_j\right) \\
			&=\sum_{i,j=1}^n(\alpha,\varepsilon_i)(\alpha,\varepsilon_j)(\varepsilon_i,\varepsilon_j)
			=\sum_{i=1}^n(\alpha,\varepsilon_i)^2.
		\end{aligned}
		\]
	\end{enumerate}
\end{solution}

\begin{problem}

	设 $\varepsilon_1,\cdots,\varepsilon_n$ 是欧氏空间 $V$ 的一组基. 证明：对任意实数
	$c_1,\cdots,c_n$，存在唯一的 $\alpha\in V$，使得
	\[
	(\varepsilon_i,\alpha)=c_i,
	\qquad i=1,2,\cdots,n.
	\]
\end{problem}
\begin{solution}
	用待定系数法. 设
	\[
	\alpha=\sum_{j=1}^n x_j\varepsilon_j.
	\]
	则 $\alpha$ 满足 $(\varepsilon_i,\alpha)=c_i$ 当且仅当其坐标 $x_1,\cdots,x_n$ 满足线性方程组
	\[
	\sum_{j=1}^n(\varepsilon_i,\varepsilon_j)x_j=c_i,
	\qquad 1\leqslant i\leqslant n.
	\]
	其系数矩阵
	\[
	G=((\varepsilon_i,\varepsilon_j))_{1\leqslant i,j\leqslant n}
	\]
	是内积关于基 $\varepsilon_1,\cdots,\varepsilon_n$ 的度量矩阵. 若 $X=(x_1,\cdots,x_n)'\ne0$，则
	\[
	X'GX=\left|\sum_{i=1}^n x_i\varepsilon_i\right|^2>0,
	\]
	所以 $G$ 可逆. 因此上述方程组有且仅有一个解，从而所求 $\alpha$ 存在且唯一.
\end{solution}

\begin{problem}
	设 $A\in M_n(\R)$ 是对角元全为正的上三角矩阵. 若 $A$ 是正交矩阵，证明 $A=E_n$.
\end{problem}
\begin{solution}
	由于 $A$ 是正交矩阵，故 $A^{-1}=A'$. 又因为 $A$ 是上三角矩阵，所以 $A^{-1}$ 也是上三角矩阵；但 $A'=A^{-1}$ 是下三角矩阵，于是 $A$ 同时为上、下三角矩阵，故
	\[
	A=\operatorname{diag}(a_1,\cdots,a_n).
	\]
	由 $AA'=E_n$ 得 $a_i^2=1$. 又 $a_i>0$，故 $a_i=1$，$1\leqslant i\leqslant n$. 因此 $A=E_n$.
\end{solution}

\begin{problem}
	证明任意可逆实 $n\times n$ 矩阵 $A$ 都可唯一分解为
	\[
	A=QR,
	\]
	其中 $Q$ 是正交矩阵，$R$ 是对角元全为正的上三角矩阵. 这称为矩阵 $A$ 的 $QR$ 分解.
\end{problem}
\begin{solution}
	记 $A$ 的列向量为 $\alpha_1,\cdots,\alpha_n$. 因 $A$ 可逆，$\alpha_1,\cdots,\alpha_n$ 线性无关. 对它们作 Schmidt 正交化并单位化，得一组标准正交向量
	\[
	\varepsilon_1,\cdots,\varepsilon_n.
	\]
	由于每个 $\varepsilon_i$ 都属于 $\Span(\alpha_1,\cdots,\alpha_i)$，且关于 $\alpha_i$ 的系数为正，故存在一个对角元全为正的上三角矩阵 $P$，使得
	\[
	(\varepsilon_1,\cdots,\varepsilon_n)=(\alpha_1,\cdots,\alpha_n)P=AP.
	\]
	令
	\[
	Q=(\varepsilon_1,\cdots,\varepsilon_n),
	\qquad R=P^{-1},
	\]
	则 $Q$ 为正交矩阵，$R$ 为对角元全为正的上三角矩阵，并且 $A=QR$.
	
	下证唯一性. 若
	\[
	A=QR=Q_1R_1,
	\]
	其中 $Q,Q_1$ 正交，$R,R_1$ 都是对角元全为正的上三角矩阵，则
	\[
	Q_1^{-1}Q=R_1R^{-1}.
	\]
	左边为正交矩阵，右边为对角元全为正的上三角矩阵. 由习题 27，二者均为 $E_n$. 故 $Q=Q_1$，进而 $R=R_1$.
\end{solution}

\begin{problem}
	\textbf{补充：Iwasawa 分解.}
	证明任意可逆实矩阵 $A$ 都可唯一写成
	\[
	A=QDN,
	\]
	其中 $Q$ 是正交矩阵，$D$ 是对角元为正实数的对角矩阵，$N$ 是对角元全为 $1$ 的上三角矩阵.
\end{problem}
\begin{solution}
	由 $QR$ 分解，$A=QR$，其中 $Q$ 正交，$R$ 是对角元全为正的上三角矩阵. 令
	\[
	D=\operatorname{diag}(r_{11},\cdots,r_{nn}),
	\qquad N=D^{-1}R.
	\]
	则 $D$ 的对角元为正，$N$ 是对角元全为 $1$ 的上三角矩阵，且 $R=DN$，从而 $A=QDN$.
	唯一性由 $QR$ 分解的唯一性以及 $R=DN$ 的显然唯一性得到.
\end{solution}

\begin{problem}

	设 $A$ 是实反对称矩阵，即 $A'=-A$. 证明 $E+A$ 可逆，且
	\[
	(E-A)(E+A)^{-1}
	\]
	是正交矩阵.
\end{problem}
\begin{solution}
	先证 $E+A$ 可逆. 若 $(E+A)X=0$，则 $AX=-X$. 于是
	\[
	X'AX=-X'X.
	\]
	但 $A'=-A$，故 $X'AX=0$，从而 $X'X=0$，即 $X=0$. 因此 $E+A$ 可逆.
	
	记
	\[
	B=(E-A)(E+A)^{-1}.
	\]
	由于 $E-A$ 与 $E+A$ 可交换，且
	\[
	B'=\left((E+A)^{-1}\right)'(E-A)'=(E-A)^{-1}(E+A),
	\]
	所以
	\[
	\begin{aligned}
		B'B
		&=(E-A)^{-1}(E+A)(E-A)(E+A)^{-1} \\
		&=E.
	\end{aligned}
	\]
	故 $B$ 是正交矩阵.
\end{solution}

\begin{problem}

	证明下列 $n\times n$ Helmert 矩阵是正交矩阵：
	\[
	H_n=
	\begin{pmatrix}
		\dfrac1{\sqrt n} & \dfrac1{\sqrt n} & \dfrac1{\sqrt n} & \dfrac1{\sqrt n} & \cdots & \dfrac1{\sqrt n} \\
		\dfrac1{\sqrt{1\cdot2}} & -\dfrac1{\sqrt{1\cdot2}} & 0 & 0 & \cdots & 0 \\
		\dfrac1{\sqrt{2\cdot3}} & \dfrac1{\sqrt{2\cdot3}} & -\dfrac2{\sqrt{2\cdot3}} & 0 & \cdots & 0 \\
		\dfrac1{\sqrt{3\cdot4}} & \dfrac1{\sqrt{3\cdot4}} & \dfrac1{\sqrt{3\cdot4}} & -\dfrac3{\sqrt{3\cdot4}} & \cdots & 0 \\
		\vdots & \vdots & \vdots & \vdots & \ddots & \vdots \\
		\dfrac1{\sqrt{(n-1)n}} & \dfrac1{\sqrt{(n-1)n}} & \dfrac1{\sqrt{(n-1)n}} & \dfrac1{\sqrt{(n-1)n}} & \cdots & -\dfrac{n-1}{\sqrt{(n-1)n}}
	\end{pmatrix}.
	\]
\end{problem}
\begin{solution}
	记 $H_n$ 的行向量为 $\alpha_0,\alpha_1,\cdots,\alpha_{n-1}$. 先有
	\[
	|\alpha_0|^2=\sum_{k=1}^n\frac1n=1.
	\]
	对 $1\leqslant i\leqslant n-1$，第 $i+1$ 行的前 $i$ 个分量为 $1/\sqrt{i(i+1)}$，第 $i+1$ 个分量为 $-i/\sqrt{i(i+1)}$，其余为 $0$. 因此
	\[
	|\alpha_i|^2
	=\sum_{k=1}^{i}\frac1{i(i+1)}+\frac{i^2}{i(i+1)}
	=\frac{i+i^2}{i(i+1)}=1.
	\]
	又对 $1\leqslant i\leqslant n-1$，
	\[
	\alpha_0\alpha_i'
	=\frac{i}{\sqrt n\sqrt{i(i+1)}}-\frac{i}{\sqrt n\sqrt{i(i+1)}}=0.
	\]
	对 $1\leqslant i<j\leqslant n-1$，
	\[
	\alpha_i\alpha_j'
	=\sum_{k=1}^{i}\frac1{\sqrt{i(i+1)}\sqrt{j(j+1)}}
	-\frac{i}{\sqrt{i(i+1)}\sqrt{j(j+1)}}=0.
	\]
	故 $H_n$ 的各行向量两两正交且长度均为 $1$，所以 $H_n$ 是正交矩阵.
\end{solution}

\begin{problem}\label{39}

	设 $V$ 是欧氏空间. 对 $\eta_1,\cdots,\eta_m\in V$，用 Schmidt 正交化作出
	\[
	\begin{aligned}
		\xi_1&=\eta_1,\\
		\xi_2&=\eta_2-\mathrm{pr}_{\xi_1}(\eta_2),\\
		&\ \vdots\\
		\xi_i&=\eta_i-\mathrm{pr}_{\xi_1}(\eta_i)-\cdots-\mathrm{pr}_{\xi_{i-1}}(\eta_i),
		\qquad 2\leqslant i\leqslant m,
	\end{aligned}
	\]
	其中
	\[
	\mathrm{pr}_{\beta}(\alpha)=
	\begin{cases}
		0, & \beta=0,\\[2mm]
		\dfrac{(\alpha,\beta)}{(\beta,\beta)}\beta, & \beta\ne0.
	\end{cases}
	\]
	记 $G(\eta_1,\cdots,\eta_m)$ 为向量组 $\eta_1,\cdots,\eta_m$ 的 Gram 矩阵. 证明：
	\begin{enumerate}
		\item $\xi_1,\xi_2,\cdots,\xi_m$ 两两正交；
		\item
		\[
		|G(\eta_1,\cdots,\eta_m)|=|G(\xi_1,\cdots,\xi_m)|=(\xi_1,\xi_1)\cdots(\xi_m,\xi_m);
		\]
		\item
		\[
		|G(\eta_1,\cdots,\eta_m)|\leqslant |\eta_1|^2\cdots |\eta_m|^2.
		\]
	\end{enumerate}
\end{problem}
\begin{solution}
	\begin{enumerate}
		\item 对 $i$ 作归纳. 设 $\xi_1,\cdots,\xi_{i-1}$ 已两两正交. 若 $1\leqslant j<i$ 且 $\xi_j\ne0$，则
		\[
		\begin{aligned}
			(\xi_i,\xi_j)
			&=(\eta_i,\xi_j)-\left(\mathrm{pr}_{\xi_j}(\eta_i),\xi_j\right) \\
			&=(\eta_i,\xi_j)-\frac{(\eta_i,\xi_j)}{(\xi_j,\xi_j)}(\xi_j,\xi_j)=0.
		\end{aligned}
		\]
		若 $\xi_j=0$，结论显然. 因此 $\xi_1,\cdots,\xi_m$ 两两正交.
		\item 由构造可知存在对角元全为 $1$ 的上三角矩阵 $T$，使得
		\[
		(\xi_1,\cdots,\xi_m)=(\eta_1,\cdots,\eta_m)T.
		\]
		于是
		\[
		G(\xi_1,\cdots,\xi_m)=T'G(\eta_1,\cdots,\eta_m)T.
		\]
		因为 $|T|=1$，所以
		\[
		|G(\xi_1,\cdots,\xi_m)|=|G(\eta_1,\cdots,\eta_m)|.
		\]
		又由 (1)，$G(\xi_1,\cdots,\xi_m)$ 是对角矩阵，故
		\[
		|G(\xi_1,\cdots,\xi_m)|=(\xi_1,\xi_1)\cdots(\xi_m,\xi_m).
		\]
		\item 由正交分解，$\eta_i$ 等于 $\xi_i$ 加上若干个与 $\xi_i$ 正交的投影项，故
		\[
		|\xi_i|\leqslant |\eta_i|,
		\qquad 1\leqslant i\leqslant m.
		\]
		结合 (2) 得
		\[
		|G(\eta_1,\cdots,\eta_m)|=\prod_{i=1}^m|\xi_i|^2
		\leqslant \prod_{i=1}^m|\eta_i|^2.
		\]
	\end{enumerate}
\end{solution}

\begin{problem}

	设 $A=(a_{ij})_{n\times n}$ 是实矩阵，记 $A$ 的列向量为 $\eta_1,\cdots,\eta_n\in\R^n$. 证明
	\[
	G(\eta_1,\cdots,\eta_n)=A'A.
	\]
	其中 $\eta_1,\cdots,\eta_n$ 的内积取标准点积. 推出 Hadamard 不等式：
	\[
	|A|^2\leqslant \prod_{j=1}^n\sum_{i=1}^n a_{ij}^2.
	\]
	等号成立当且仅当
	\[
	\sum_{k=1}^n a_{ki}a_{kj}=0,
	\qquad i\ne j,
	\]
	即列向量两两正交，或 $A$ 有零列.
\end{problem}
\begin{solution}
	由于 $A=(\eta_1,\cdots,\eta_n)$，故
	\[
	A'A=\bigl((\eta_i,\eta_j)\bigr)_{1\leqslant i,j\leqslant n}=G(\eta_1,\cdots,\eta_n).
	\]
	于是由\ref{39}，
	\[
	|A|^2=|A'A|=|G(\eta_1,\cdots,\eta_n)|
	\leqslant \prod_{j=1}^n |\eta_j|^2
	=\prod_{j=1}^n\sum_{i=1}^n a_{ij}^2.
	\]
	等号成立正是 Gram 行列式不等式的等号情形，即列向量 $\eta_1,
	\cdots,\eta_n$ 两两正交，或者某个 $\eta_j=0$；写成坐标就是
	\[
	(\eta_i,\eta_j)=\sum_{k=1}^n a_{ki}a_{kj}=0,
	\qquad i\ne j,
	\]
	或 $A$ 有零列.
\end{solution}

\begin{problem}

	设 $U,W$ 是有限维欧氏空间 $V$ 的子空间. 证明：
	\begin{enumerate}
		\item $W=(W^\perp)^\perp$；
		\item 若 $U\subseteq W$，则 $W^\perp\subseteq U^\perp$；
		\item $(U+W)^\perp=U^\perp\cap W^\perp$；
		\item $(U\cap W)^\perp=U^\perp+W^\perp$.
	\end{enumerate}
\end{problem}
\begin{solution}
	\begin{enumerate}
		\item 显然 $W\subseteq (W^\perp)^\perp$. 又
		\[
		\dim (W^\perp)^\perp=\dim V-\dim W^\perp=\dim W,
		\]
		故 $W=(W^\perp)^\perp$.
		\item 若 $\alpha\in W^\perp$，则对任意 $\beta\in U\subseteq W$，有 $(\alpha,\beta)=0$，故 $\alpha\in U^\perp$.
		\item 对任意 $\alpha\in V$，
		\[
		\begin{aligned}
			\alpha\in (U+W)^\perp
			&\Longleftrightarrow (\alpha,u+w)=0,
			\quad \forall u\in U,\ w\in W \\
			&\Longleftrightarrow \alpha\in U^\perp\text{ 且 }\alpha\in W^\perp \\
			&\Longleftrightarrow \alpha\in U^\perp\cap W^\perp.
		\end{aligned}
		\]
		\item 由 (3) 作用于 $U^\perp$ 与 $W^\perp$，并用 (1)，得
		\[
		\begin{aligned}
			(U^\perp+W^\perp)^\perp
			&=(U^\perp)^\perp\cap(W^\perp)^\perp \\
			&=U\cap W.
		\end{aligned}
		\]
		两边再取正交补，得 $(U\cap W)^\perp=U^\perp+W^\perp$.
	\end{enumerate}
\end{solution}

\begin{problem}

	设 $\beta_1,\cdots,\beta_m$ 是 $V$ 中一组两两正交的非零向量. 证明对任意 $\beta\in V$，有
	\[
	\sum_{i=1}^m\frac{(\beta,\beta_i)^2}{(\beta_i,\beta_i)}\leqslant (\beta,\beta).
	\]
\end{problem}
\begin{solution}
	令
	\[
	p=\sum_{i=1}^m\frac{(\beta,\beta_i)}{(\beta_i,\beta_i)}\beta_i.
	\]
	则对每个 $i$，有 $(\beta-p,\beta_i)=0$，故 $\beta-p\perp p$. 因而
	\[
	|\beta|^2=|p|^2+|\beta-p|^2\geqslant |p|^2.
	\]
	又由 $\beta_1,\cdots,\beta_m$ 两两正交，
	\[
	|p|^2=\sum_{i=1}^m\frac{(\beta,\beta_i)^2}{(\beta_i,\beta_i)}.
	\]
	于是得到所证不等式. 这正是 Bessel 不等式.
\end{solution}

\begin{problem}

	设 $W$ 是欧氏空间 $V$ 的子空间. 定义向量 $\alpha$ 到 $W$ 的距离为
	\[
	d(\alpha,W)=\min\{\,|\alpha-\beta|\mid \beta\in W\,\}.
	\]
	证明
	\[
	d(\alpha,W)=|\alpha-\alpha_1|,
	\]
	其中 $\alpha_1$ 是 $\alpha$ 在 $W$ 上的正交投影. 并给出 $V=\R^3$ 时的几何解释.
\end{problem}
\begin{solution}
	设
	\[
	\alpha=\alpha_1+\alpha_2,
	\qquad \alpha_1\in W,
	\quad \alpha_2\in W^\perp.
	\]
	对任意 $\beta\in W$，有 $\alpha_1-\beta\in W$，故 $\alpha_1-\beta\perp \alpha_2$. 因此
	\[
	\begin{aligned}
		|\alpha-\beta|^2
		&=|(\alpha_1-\beta)+\alpha_2|^2 \\
		&=|\alpha_1-\beta|^2+|\alpha_2|^2
		\geqslant |\alpha_2|^2=|\alpha-\alpha_1|^2.
	\end{aligned}
	\]
	当 $\beta=\alpha_1$ 时等号成立，所以
	\[
	d(\alpha,W)=|\alpha-\alpha_1|.
	\]
	在 $\R^3$ 中，若 $W$ 是一条直线或一个平面，则该距离就是点 $\alpha$ 到这条直线或这个平面的垂线段长度.
\end{solution}

\begin{problem}

	设 $W$ 是欧氏空间 $V$ 的子空间. 定义向量 $\alpha$ 与 $W$ 的夹角为
	\[
	\angle(\alpha,W)=\min\{\,\angle(\alpha,\beta)\mid 0\ne\beta\in W\,\}.
	\]
	证明
	\[
	\angle(\alpha,W)=\angle(\alpha,\alpha_1),
	\]
	其中 $\alpha_1$ 是 $\alpha$ 在 $W$ 上的正交投影. 并给出 $V=\R^3$ 时的几何解释.
\end{problem}
\begin{solution}
	只需证明在 $W$ 中，$\alpha_1$ 的方向使投影长度最大. 设 $0\ne\beta\in W$. 因为 $\alpha-\alpha_1\in W^\perp$，而 $\mathrm{pr}_{\beta}(\alpha)\in W$，所以
	\[
	(\alpha-\alpha_1,\mathrm{pr}_{\beta}(\alpha))=0.
	\]
	另一方面，由投影定义，
	\[
	(\alpha-\mathrm{pr}_{\beta}(\alpha),\mathrm{pr}_{\beta}(\alpha))=0.
	\]
	两式相减得
	\[
	(\alpha_1-\mathrm{pr}_{\beta}(\alpha),\mathrm{pr}_{\beta}(\alpha))=0.
	\]
	于是
	\[
	|\alpha_1|^2=|\alpha_1-\mathrm{pr}_{\beta}(\alpha)|^2+|\mathrm{pr}_{\beta}(\alpha)|^2
	\geqslant |\mathrm{pr}_{\beta}(\alpha)|^2.
	\]
	又 $\mathrm{pr}_{\alpha_1}(\alpha)=\alpha_1$，故 $\alpha_1$ 给出的投影长度最大. 由于余弦函数在 $[0,\pi/2]$ 上单调递减，故
	\[
	\angle(\alpha,W)=\angle(\alpha,\alpha_1).
	\]
	在 $\R^3$ 中，若 $W$ 是一个平面，则 $\angle(\alpha,W)$ 就是直线 $\alpha$ 与平面的夹角.
\end{solution}

\begin{problem}

	设 $\mathcal A:\R^3\to\R^3$ 是正交变换.
	\begin{enumerate}
		\item 若 $\det\mathcal A=1$，则 $\mathcal A$ 是绕过原点的一条轴的旋转，该轴由属于特征值 $1$ 的特征向量生成；
		\item 若 $\det\mathcal A=-1$，则 $\mathcal A$ 是关于某平面 $P$ 的反射，再复合绕垂直于 $P$ 的直线 $\ell$ 的旋转；直线 $\ell$ 由属于特征值 $-1$ 的特征向量生成；
		\item 对第二类正交变换
		\[
		\mathcal A(x,y,z)=(-x,y,z),
		\]
		求出 (2) 中的平面 $P$ 与直线 $\ell$.
	\end{enumerate}
\end{problem}
\begin{solution}
	三维实线性空间上的任意线性变换都有实特征值，而正交变换的实特征值只能为 $\pm1$.
	\begin{enumerate}
		\item 若 $\det\mathcal A=1$，则 $\mathcal A$ 必有特征值 $1$. 取单位特征向量 $u$，令 $\ell=\Span(u)$. 由于 $\mathcal A$ 正交，$\ell^\perp$ 是不变子空间，且 $\mathcal A|_{\ell^\perp}$ 是平面上的正交变换. 其行列式仍为 $1$，故为平面旋转. 因此 $\mathcal A$ 是绕 $\ell$ 的旋转.
		\item 若 $\det\mathcal A=-1$，则 $\mathcal A$ 必有特征值 $-1$. 取单位特征向量 $u$，令 $\ell=\Span(u)$，$P=\ell^\perp$. 则 $P$ 是不变平面，且 $\mathcal A|_P$ 是平面旋转. 设 $S_P$ 为关于平面 $P$ 的反射，则 $S_P$ 在 $P$ 上恒等，在 $\ell$ 上变号. 因而 $\mathcal A$ 可看作先作 $S_P$，再绕 $\ell$ 作平面 $P$ 上对应的旋转.
		\item 对 $\mathcal A(x,y,z)=(-x,y,z)$，属于特征值 $-1$ 的特征向量可取 $(1,0,0)$. 故
		\[
		\ell=\Span\{(1,0,0)\},
		\qquad
		P=\ell^\perp=\{(0,y,z):y,z\in\R\}.
		\]
		此时绕 $\ell$ 的旋转就是恒等旋转，$\mathcal A$ 本身即关于 $P$ 的反射.
	\end{enumerate}
\end{solution}

\begin{problem}
	设 $V$ 是欧氏空间. 设映射 $\mathcal A:V\to V$ 满足
	\[
	|\alpha-\beta|=|\mathcal A(\alpha)-\mathcal A(\beta)|,
	\qquad \alpha,\beta\in V.
	\]
	证明存在正交变换 $\mathcal B$ 与 $\gamma\in V$，使得
	\[
	\mathcal A(\alpha)=\mathcal B\alpha+\gamma,
	\qquad \alpha\in V.
	\]
\end{problem}
\begin{solution}
	令
	\[
	\gamma=\mathcal A(0),
	\qquad
	\mathcal B(\alpha)=\mathcal A(\alpha)-\gamma.
	\]
	则 $\mathcal B(0)=0$，并且
	\[
	|\mathcal B(\alpha)-\mathcal B(\beta)|=|\alpha-\beta|.
	\]
	特别地，$|\mathcal B(\alpha)|=|\alpha|$. 由极化恒等式，对任意 $\alpha,\beta\in V$，
	\[
	\begin{aligned}
		(\mathcal B\alpha,\mathcal B\beta)
		&=\frac12\left(|\mathcal B\alpha|^2+|\mathcal B\beta|^2-|\mathcal B\alpha-\mathcal B\beta|^2\right) \\
		&=\frac12\left(|\alpha|^2+|\beta|^2-|\alpha-\beta|^2\right)
		=(\alpha,\beta).
	\end{aligned}
	\]
	由此展开可得
	\[
	|\mathcal B(\alpha+\beta)-\mathcal B\alpha-\mathcal B\beta|^2=0,
	\qquad
	|\mathcal B(c\alpha)-c\mathcal B\alpha|^2=0,
	\]
	故 $\mathcal B$ 是线性映射. 又它保持内积，所以 $\mathcal B$ 是正交变换. 因而
	\[
	\mathcal A(\alpha)=\mathcal B\alpha+\gamma.
	\]
\end{solution}

\begin{problem}
	证明实反对称矩阵的所有特征值都是纯虚数或 $0$.
\end{problem}
\begin{solution}
	设 $A$ 是 $n$ 阶实反对称矩阵，$\lambda_0$ 是 $A$ 的任意复特征值. 取非零向量
	$\xi=(z_1,\cdots,z_n)'\in\C^n$，使得
	\[
	A\xi=\lambda_0\xi.
	\]
	由于 $A$ 是实矩阵，取复共轭得
	\[
	A\overline{\xi}=\overline{\lambda_0}\,\overline{\xi}.
	\]
	又 $A'=-A$，于是转置上式得
	\[
	(A\overline{\xi})'=\overline{\xi}^{\,\prime}A'=-\overline{\xi}^{\,\prime}A
	=\overline{\lambda_0}\,\overline{\xi}^{\,\prime}.
	\]
	右乘 $\xi$，得
	\[
	-\overline{\xi}^{\,\prime}A\xi
	=-\lambda_0\overline{\xi}^{\,\prime}\xi
	=\overline{\lambda_0}\,\overline{\xi}^{\,\prime}\xi.
	\]
	即
	\[
	(\lambda_0+\overline{\lambda_0})\overline{\xi}^{\,\prime}\xi=0.
	\]
	但
	\[
	\overline{\xi}^{\,\prime}\xi=|z_1|^2+\cdots+|z_n|^2>0,
	\]
	所以 $\lambda_0+\overline{\lambda_0}=0$. 因此 $\lambda_0$ 为纯虚数或 $0$.
\end{solution}

\begin{problem}
	设 $\mathcal A$ 是欧氏空间 $V$ 上的对称变换. 证明
	\[
	\Im(\mathcal A)=(\Ker\mathcal A)^\perp,
	\qquad
	\Ker\mathcal A=(\Im \mathcal A)^\perp.
	\]
	特别地，
	\[
	V=\Im(\mathcal A)\oplus\Ker\mathcal A.
	\]
\end{problem}
\begin{solution}
	利用对称变换可对角化. 取 $V$ 的一组标准正交基
	\[
	e_1,\cdots,e_r,e_{r+1},\cdots,e_n,
	\]
	使得 $\mathcal A$ 在该基下的矩阵为
	\[
	\operatorname{diag}(a_1,\cdots,a_r,0,\cdots,0),
	\qquad a_i\ne0.
	\]
	于是
	\[
	\Im(\mathcal A)=\Span(e_1,\cdots,e_r),
	\qquad
	\Ker\mathcal A=\Span(e_{r+1},\cdots,e_n).
	\]
	这两个子空间互为正交补，故
	\[
	\Im(\mathcal A)=(\Ker\mathcal A)^\perp,
	\qquad
	\Ker\mathcal A=(\Im \mathcal A)^\perp.
	\]
	于是
	\[
	V=\Im(\mathcal A)\oplus\Ker\mathcal A.
	\]
\end{solution}

\begin{problem}
	设 $A$ 是 $n\times n$ 实对称矩阵，其互异特征值为 $\lambda_1,\lambda_2,\cdots,\lambda_r$. 证明存在矩阵 $A_1,A_2,\cdots,A_r$，使得
	\begin{enumerate}
		\item $A_1+A_2+\cdots+A_r=E_n$；
		\item $A_i^2=A_i$，$1\leqslant i\leqslant r$；
		\item $A_iA_j=0$，$i\ne j$，$1\leqslant i,j\leqslant r$；
		\item $A=\lambda_1A_1+\lambda_2A_2+\cdots+\lambda_rA_r$.
	\end{enumerate}
\end{problem}
\begin{solution}
	因为 $A$ 是实对称矩阵，所以存在正交矩阵 $T$，使得
	\[
	T^{-1}AT=\operatorname{diag}(\lambda_1E_{t_1},\lambda_2E_{t_2},\cdots,\lambda_rE_{t_r}),
	\]
	其中 $t_i$ 是特征值 $\lambda_i$ 的重数. 令
	\[
	F_i=
	\begin{pmatrix}
		0 &        &        &        & 0\\
		& \ddots &        &        &  \\
		&        & E_{t_i}&        &  \\
		&        &        & \ddots &  \\
		0 &        &        &        & 0
	\end{pmatrix},
	\qquad
	A_i=TF_iT^{-1},
	\quad 1\leqslant i\leqslant r,
	\]
	其中 $F_i$ 是与第 $i$ 个特征值对应的分块单位矩阵，其余分块为零. 于是
	\[
	F_1+\cdots+F_r=E_n,
	\qquad F_i^2=F_i,
	\qquad F_iF_j=0\quad(i\ne j).
	\]
	从而
	\[
	A_1+\cdots+A_r=T(E_n)T^{-1}=E_n,
	\]
	并且
	\[
	A_i^2=TF_i^2T^{-1}=TF_iT^{-1}=A_i,
	\qquad
	A_iA_j=TF_iF_jT^{-1}=0\quad(i\ne j).
	\]
	最后，
	\[
	\begin{aligned}
		\lambda_1A_1+\cdots+\lambda_rA_r
		&=T(\lambda_1F_1+\cdots+\lambda_rF_r)T^{-1} \\
		&=T\operatorname{diag}(\lambda_1E_{t_1},\cdots,\lambda_rE_{t_r})T^{-1}=A.
	\end{aligned}
	\]
\end{solution}
